\documentclass[a4paper,11pt]{article}
\pdfoutput=1

\usepackage[table]{xcolor}
\usepackage{lmodern}
\usepackage{jcappub}
\usepackage[toc,page]{appendix}
\usepackage{subcaption}
\usepackage{mathtools}
\usepackage{orcidlink}
\usepackage{multirow}
\usepackage[T1]{fontenc}
\usepackage{float}
\usepackage{bm}
\usepackage{tcolorbox}

\tcbuselibrary{skins}
\tcbuselibrary{breakable}

\newcommand{\semibold}[1]{{\fontseries{b}\selectfont{#1}}}
\newcommand{\sansbold}[1]{{\sffamily\fontseries{sbc}\selectfont{#1}}}

\newcommand{\para}[1]{\par\vspace{2mm}\noindent\semibold{{#1.}---}\ignorespaces}

\DeclareMathOperator{\Or}{O}
\DeclareMathOperator{\tr}{tr}
\DeclareMathOperator{\NewRe}{Re}
\DeclareMathOperator{\NewIm}{Im}
\renewcommand{\Re}{\NewRe}
\renewcommand{\Im}{\NewIm}

\newcommand{\transpose}{\mathsf{T}}

\renewcommand{\leq}{\leqslant}
\renewcommand{\geq}{\geqslant}

\definecolor{SussexCobaltBlue}{HTML}{1d4289}
\definecolor{SussexDeepAquamarine}{HTML}{007a78}
\definecolor{SussexPowderBlue}{HTML}{7da1c4}
\definecolor{SussexCornYellow}{HTML}{f2c75c}
\definecolor{SussexChinaRose}{HTML}{be84a3}
\definecolor{SussexBurntOrange}{HTML}{dc582a}

\hypersetup{colorlinks=true,citecolor=SussexDeepAquamarine,linkcolor=SussexCobaltBlue,urlcolor=SussexBurntOrange}

\newcommand{\etal}{\emph{et al.}}

\renewcommand{\d}{\mathrm{d}}
\newcommand{\e}[1]{\mathrm{e}^{{#1}}}
\newcommand{\im}{\mathrm{i}}
\newcommand{\Dm}[1]{\mathcal{D}#1}            % functional-integral measure, e.g. \Dm\hat\field
\newcommand{\Mp}{M_{\mathrm{P}}}
\newcommand{\vect}[1]{\bm{\mathrm{{#1}}}}

\DeclareMathOperator{\grad}{\nabla}
\DeclareMathOperator{\sinc}{sinc}

\newcommand{\GeV}{\text{GeV}}

% spread out arrays / text
\renewcommand{\arraystretch}{1.15}
\renewcommand{\baselinestretch}{1.1}

\newenvironment{panel}{\begin{tcolorbox}[enhanced,breakable,colback=gray!10,colframe=gray!30]\small}{\end{tcolorbox}}
\newenvironment{todo}{\begin{tcolorbox}[enhanced,breakable,colback=SussexBurntOrange!10,colframe=SussexBurntOrange!40]\small{\textsf{\textbf{Open issue:} }}}{\end{tcolorbox}}
\newenvironment{calculation}{\begin{tcolorbox}[enhanced,breakable,colback=SussexPowderBlue!10,colframe=SussexCobaltBlue!30]\small}{\end{tcolorbox}}

%
% ----------------------------------------------------------------------
% Notation macros. Collected here so that definitions can be changed
% in one place. Grouped by role.
% ----------------------------------------------------------------------
%

% -- e-folding time and background quantities --
\newcommand{\Ny}{N}                          % e-fold "time" variable
\newcommand{\Ninit}{N_{\mathrm{init}}}       % start of instanton transition
\newcommand{\Nfinal}{N_{\mathrm{final}}}     % end of instanton transition
\newcommand{\DN}{\Delta N}                   % duration of transition, N_init - N_final on background
\newcommand{\dNstar}{\delta N_\star}         % instanton excess e-folds
\newcommand{\epsSR}{\epsilon_1}              % first slow-roll parameter
\newcommand{\Hub}{H}                         % Hubble rate
\newcommand{\sfactor}{a}                     % scale factor
\newcommand{\atloc}{(\field,\mom)}           % shorthand: evaluated at the local phase-space
                                              % point (field,mom); distinguishes the LOCAL,
                                              % dynamical use of H, eps_1 (Friedmann-equation
                                              % functions of the local trajectory, exactly as
                                              % in AbstractPotential.H_sq/epsilon) from the
                                              % GLOBAL background aH(N) used only to define
                                              % the coordinate y and the gradient-term prefactor

% -- field content --
\newcommand{\field}{\phi}                    % inflaton field
\newcommand{\mom}{\pi}                       % conjugate momentum, d(field)/dN
\newcommand{\Vpot}{V}                        % potential
\newcommand{\phinl}{\field_{\mathrm{nl}}}    % noiseless background trajectory
\newcommand{\pinl}{\mom_{\mathrm{nl}}}       % noiseless background momentum
\newcommand{\phiinit}{\field_{\mathrm{init}}}
\newcommand{\phiend}{\field_{\mathrm{end}}}

% -- noise / MSR response fields --
\newcommand{\noise}[1]{\xi_{#1}}             % generic noise, e.g. \noise{\field}
\newcommand{\rfield}{\widetilde\field}       % Real saddle-point response field conjugate to
                                              % \field. NOT literally the naive delta-functional
                                              % auxiliary field \hat\field from the MSR
                                              % construction: related by \hat\field=i\rfield
                                              % (Section~\ref{sec:msr-action}); \rfield is the
                                              % field rotated onto the imaginary axis, real by
                                              % construction, and is what appears in the action,
                                              % the instanton equations, and the code (P1).
\newcommand{\rmom}{\widetilde\mom}           % Real saddle-point response field conjugate to
                                              % \mom, with the same caveat: \hat\mom=i\rmom.
                                              % Corresponds to P2 in the code.

% -- shell / layer bookkeeping --
\newcommand{\Nlay}[1]{N_{#1}}                % comoving population of shell i, N_i
\newcommand{\rad}{r}                         % comoving radial coordinate

% -- the onion coordinate y and its complement u --
\newcommand{\yco}{y}                         % onion coordinate, y=-1 fixed outer edge,
                                              % y=+1 current horizon (exterior-only, log-radial)
\newcommand{\uco}{u}                         % u = 1-y (legacy; unused in current coordinate)
\newcommand{\rout}{\rad_{\mathrm{out}}}      % comoving radius of outer edge of overdensity
\newcommand{\rphys}{\rad_{\mathrm{phys}}}    % present-day physical scale (Leach-Liddle), distinct
                                              % from comoving \rad and areal R (Section~\ref{sec:scale-assignment})
\newcommand{\aHo}{(\sfactor \Hub)_0}         % aH evaluated at N_init
\newcommand{\aH}[1][]{(\sfactor \Hub)#1}     % aH, optionally with an argument e.g. \aH{(N)}
\newcommand{\aHloc}{\sfactor(\Ny)\Hub\atloc} % LOCAL comoving-Hubble-rate quantity for the
                                              % gradient (Laplacian) coupling term: a(N) is a
                                              % pure function of the grid N (N is itself each
                                              % shell's own e-fold clock, so a(N)=a_0 e^{N-Ninit}
                                              % is exact, not tied to any trajectory), while H
                                              % is local, H(field,mom). Contrast \aH[(\Ny)],
                                              % retained for the coordinate/geometric roles
                                              % (defining y, aHo, n(y,N), k_sigma) where a single
                                              % shared comoving map is required.
\newcommand{\rH}{\rad_H}                     % current horizon radius, evaluated at the CORE's
                                              % own local field values: rH(N) = 1/[a(N)H(1,N)]
\newcommand{\Deltas}{\Delta s}                % Deltas(N) = ln(rout/rH(N)): log-radial extent
                                              % of the resolved (exterior) domain
\newcommand{\epscore}{\epsSR^{\mathrm{core}}} % epsilon_1 evaluated at the core, y=+1
\newcommand{\measure}{\mu}                   % self-adjoint measure for Lop(y,N) up to C(N): exponential
                                              % in y, N-dependent through Deltas(N)
\newcommand{\advcoef}{A}                     % advection coefficient in (y,N) coordinates,
                                              % nonzero for the log-radial horizon-excising map
\newcommand{\weightfn}{w}                    % w(y) = mu(y,N)*p(y), where p=e^{Deltas y} is the
                                              % coefficient of Q'' in Lop; equivalently mu = w^3.
                                              % Used in the self-adjointness check (Section~\ref{sec:msr-action})

% -- LEGACY: sinc-eigenfunction spectral expansion, superseded by the
% collocation scheme (Section~\ref{sec:collocation-scheme}). Retained,
% unused, only because a handful of macro names below are reused for
% collocation-grid quantities with a different meaning; do not rely on
% the "mode coefficient" comments below, which describe the old scheme.
\newcommand{\Lop}{\mathcal{L}_\yco}          % radial mean-field Laplacian in y (now the
                                              % log-radial, horizon-excising operator of
                                              % Section~\ref{sec:coordinate}, not the old sinc one)
\newcommand{\modeidx}{n}                     % mode index
\newcommand{\basisfn}{f_\modeidx}            % sinc basis function
\newcommand{\basisfnx}[1]{f_{#1}}            % sinc basis function, explicit index
\newcommand{\eigval}{\kappa_\modeidx}        % eigenvalue of \Lop, = -(n pi)^2
\newcommand{\modenorm}{N_\modeidx}           % norm of basis function f_n in the (1-y)^2 dy measure

% -- mode coefficients --
\newcommand{\coefa}{a_\modeidx}              % mode coefficient of (field - noiseless)
\newcommand{\coefax}[1]{a_{#1}}              % mode coefficient of (field - noiseless), explicit index
\newcommand{\coefb}{b_\modeidx}              % mode coefficient of (mom - noiseless)
\newcommand{\rcoefa}{\widetilde{a}_\modeidx} % mode coefficient of response field \rfield
\newcommand{\rcoefax}[1]{\widetilde{a}_{#1}} % mode coefficient of \rfield, explicit index
\newcommand{\rcoefb}{\widetilde{b}_\modeidx} % mode coefficient of response field \rmom

\newcommand{\rcoefbx}[1]{\widetilde{b}_{#1}} % mode coefficient of \rmom, explicit index
\newcommand{\projone}{[1]_\modeidx}          % projection of the constant function 1 onto
                                              % mode n in the u^2 du measure; = 2(-1)^{n+1}
% -- noise amplitude / covariance functions --
\newcommand{\Dnoise}[1]{D_{#1}}              % D_field(y,N), D_mom(y,N), D_{field mom}(y,N):
                                              % shell-averaged noise variance/covariance
                                              % densities (the "2D" convention: matches the
                                              % factor of 2 multiplying D11,D12,D22 in the
                                              % existing FullInstanton forward-pass equations)
\newcommand{\Dcm}[1]{\mathsf{D}_{#1}}        % single-patch (undiluted) diffusion matrix
                                              % element in the codebase convention, i.e.
                                              % DiffusionModel.D_matrix(phi,pi) -> (D11,D12,D22);
                                              % \Pspec{fg} = 2\,\Dcm{fg}\atloc
\newcommand{\Pspec}[1]{P_{#1}}               % reduced power spectra of Briaud et al.
\newcommand{\ncount}{n}                      % Hubble-volume count per unit shell, n(y,N)

% -- volume element, noise amplitude, Hamiltonian (added October 2026) --
\newcommand{\volel}{\mathcal{V}}             % physical volume element per unit y:
                                              % V(y,N) = dV/dy = 2 pi rout^3 Deltas e^{-(3/2)(y+1) Deltas}
\newcommand{\Cpref}{\mathcal{C}}             % y-independent prefactor, V = C(N) mu(y,N)
\newcommand{\Sigvol}[1]{\varSigma_{#1}}      % noise amplitude in the V-convention:
                                              % Sigma_ij = V D_ij = 2 Dcm_ij v_H
\newcommand{\vH}{v_H}                        % comoving volume of one Hubble patch, (4 pi/3)/(aH_loc)^3
\newcommand{\Hfp}{H_{\mathrm{FP}}}            % Fokker-Planck (Freidlin-Wentzell) Hamiltonian
\newcommand{\Hh}{H_h}                         % its discrete (collocation) version
\newcommand{\pfield}{p_\field}                % canonical momentum conjugate to phi, = V rfield
\newcommand{\pmom}{p_\mom}                    % canonical momentum conjugate to pi, = V rmom
\newcommand{\wspeed}{c}                       % characteristic (wave) speed in y
\newcommand{\win}{w_{\mathrm{in}}}            % incoming Riemann invariant at the core
\newcommand{\Fcore}{\mathcal{F}}              % non-advective core forcing, superscripts phi/pi
\newcommand{\Kcoef}{K}                        % denominator coefficient in the chain-rule Deltas-dot
\newcommand{\rHbg}{\rad_H^{\mathrm{bg}}}      % background (noiseless) comoving horizon
\newcommand{\Hen}{E}                          % H-energy (control functional), not a physical energy
\newcommand{\Hnorm}{\mathsf{H}}               % SBP norm matrix
\newcommand{\gpref}{g}                        % gradient prefactor g = e^{-2 Deltas_loc}
\newcommand{\core}{\mathrm{c}}                % subscript: evaluated at the core node y=+1
\newcommand{\Ntotal}{N_{\mathrm{total}}}      % duration of the transition, Delta N + delta N_star

% -- misc --
\newcommand{\lag}{\lambda}                   % Lagrange multiplier for core endpoint condition
\newcommand{\zprofile}{\zeta}                % curvature perturbation profile

\title{The onion model: a mean-field gradient-corrected instanton\\
for the spacetime embedding of PBH-forming trajectories}

\author{David Seery}
\affiliation{Astronomy Centre, University of Sussex, Brighton BN1 9QH, United Kingdom}
\emailAdd{D.Seery@sussex.ac.uk}

\begin{document}
\maketitle

% \begin{abstract}
% \noindent
% These notes record the derivation of a mean-field, gradient-corrected
% extension of the instanton method used elsewhere in this project to compute
% primordial black hole (PBH) formation probabilities in stochastic inflation.
% In the existing prescription (`\texttt{FullInstanton}'/`\texttt{SlowRollInstanton}'),
% a single instanton trajectory is assumed to describe the assembly of the
% central region of a spherically symmetric overdensity, while surrounding
% Hubble volumes are assumed to detach and evolve noiselessly once they exit
% the active core. This neglects the gradient coupling between the detaching
% shell and the central region identified by Briaud \etal\ in their treatment
% of gradient corrections to stochastic inflation. Here we build a mean-field
% `onion model': angle-averaged Langevin equations for the field and momentum
% in concentric comoving shells, including the leading gradient coupling
% between shells, converted to a Martin--Siggia--Rose (MSR) representation,
% with the gradient-corrected instanton found as the saddle point of a global
% boundary value problem. We derive the appropriate comoving coordinate for
% the shell structure, construct the MSR action with the correct volume
% measure, derive the instanton (Euler--Lagrange) equations and their
% boundary conditions, construct an analytic eigenbasis for the spatial
% operator, and set out a spectral collocation scheme for the numerical
% solution. Two open theoretical issues are flagged explicitly and are not
% yet resolved: the correct probabilistic interpretation of the two-dimensional
% MSR action relative to the existing single-trajectory instanton, and the
% precise sense in which the terminal (`largest time') boundary condition on
% the MSR response fields can be placed at an arbitrary time later than the
% end of the transition.
% \end{abstract}

\tableofcontents

% ======================================================================
\section{Motivation}
\label{sec:motivation}
% ======================================================================

In the existing instanton approach to PBH formation, a single trajectory
$\field(\Ny)$, $\mom(\Ny) = \d\field/\d\Ny$ is found as the saddle point of
the Martin--Siggia--Rose (MSR) action for the central Hubble patch, subject
to boundary conditions fixing the initial field value $\phiinit$, the final
field value $\phiend$, and the duration $\dNstar$ of the anomalous
(noise-driven) excursion relative to the noiseless background. The
resulting curvature perturbation profile $\zprofile(\rad)$ is built by
assuming that, at each step of the assembly, the current Hubble volume
splits into a central region --- which continues to be described by the
instanton --- and a surrounding shell, which detaches and thereafter
evolves \emph{noiselessly} along the background trajectory.

This is a reasonable first approximation, but it neglects the gradient
coupling between the detaching shell and the central region. Briaud
\etal\ have shown that adjacent Hubble volumes do not decouple completely
upon separation: the noise driving neighbouring patches is correlated on
the coarse-graining scale, and this correlation sources a residual gradient
interaction that persists for of order one $e$-fold after separation. This
coupling makes it more costly, in the sense of MSR action, to build a
sharply peaked (`Cerro Torre') profile than the noiseless-detachment
prescription would suggest, because the detaching shell resists being left
behind by the central assembly.

\begin{panel}
The Briaud \etal\ formalism itself is not directly applicable to the
rare-event geometry we are interested in, because it is built for
statistically typical patches relaxing towards the background, whereas here
the central patch is conditioned to be an extreme fluctuation. The
construction below is therefore not simply an import of their results:
we use their identification of the leading gradient coupling and noise
correlation structure, but the conditioning is handled correctly by
extremizing a single joint action for the whole configuration, rather than
by substituting an unconditional noise correlation function into a Langevin
equation by hand.
\end{panel}

The goal of the `onion model' developed here is to replace the ad hoc
noiseless-detachment rule with a dynamical boundary condition derived from
an action principle, so that the cost of assembling a sharp profile ---
and the associated smoothing of the profile by gradient coupling to the
surrounding shells --- emerges from the saddle point rather than being
imposed by hand.

% ======================================================================
\section{Single-patch Langevin system and shell averaging}
\label{sec:shell-averaging}
% ======================================================================

\subsection{The un-coarse-grained system}

Working to the order in the gradient expansion retained by Briaud \etal,
the Langevin equations for the coarse-grained field at a comoving point
$\vect{x}$ are
\begin{subequations}
\label{eq:patch-langevin}
\begin{align}
    \frac{\d \field}{\d \Ny}(\vect{x}, \Ny)
    & =
    \mom(\vect{x}, \Ny) + \noise{\field}(\vect{x}, \Ny) ,
    \\
    \frac{\d \mom}{\d \Ny}(\vect{x}, \Ny)
    & =
    -(3-\epsSR) \mom(\vect{x},\Ny)
    - \frac{\Vpot'(\field)}{\Hub^2}
    + \frac{1}{\aHloc^2} \grad^2 \field(\vect{x}, \Ny)
    + \noise{\mom}(\vect{x}, \Ny) ,
\end{align}
\end{subequations}
where $\grad^2$ is the comoving spatial Laplacian and the noises have
covariance
\begin{equation}
    \big\langle \noise{f}(\vect{x}, \Ny) \, \noise{g}(\vect{y}, \Ny') \big\rangle
    =
    \Pspec{fg}(\Ny) \, \delta_D(\Ny - \Ny') \, (1-\epsSR) \,
    \sinc\!\big( k_\sigma |\vect{x}-\vect{y}| \big) ,
    \qquad f, g \in \{\field, \mom\} ,
    \label{eq:noise-covariance}
\end{equation}
with $k_\sigma = \sigma \aH$ the coarse-graining scale ($\sigma \ll 1$) and
$\Pspec{fg}$ the reduced power spectra. The noise is white in $\Ny$ and
correlated between comoving points only within the coarse-graining radius
$k_\sigma^{-1}$.

\subsection{Shell averaging}
\begin{panel}
\semibold{$H$, $\epsSR$, $\Pspec{fg}$ here are local, not background,
quantities -- and so is $H$ inside the gradient-term coefficient.}
Throughout \eqref{eq:patch-langevin}--\eqref{eq:noise-covariance}
$\Hub$, $\epsSR$, and $\Pspec{fg}$ should be read as the usual
Friedmann-equation functions of the local phase-space point
$\field(\vect{x},\Ny),\mom(\vect{x},\Ny)$ -- exactly as in the single-patch
stochastic-$\delta N$ formalism this section is built on, and as
\texttt{AbstractPotential.H\_sq}/\texttt{epsilon} and
\texttt{DiffusionModel.D\_matrix} already treat them in the codebase --
\emph{not} as fixed background values. This is unremarkable at the
single-patch level, where there is only one field value in play, but it
matters once shells with genuinely different field content are compared
below and, later, once the shell/mode structure is built in
$(\yco,\Ny)$: see the explicit local-evaluation prescription in
Section~\ref{sec:noise-y} and the corrected equations of
Section~\ref{sec:instanton-eqs}, which this un-coarse-grained system should
be understood to feed into consistently. The same applies to the $\aH$
appearing in the gradient term $\aHloc^{-2}\grad^2\field$: $a(\Ny)$ is a
pure function of the grid $\Ny$ (exact, not an approximation, since $\Ny$
is by construction each patch's own e-fold clock -- $a(\Ny)=a_0\e{\Ny-\Ninit}$
regardless of trajectory), but $H$ there is the same local
$H\atloc$, not a value pinned to any reference trajectory. This is
distinct from the genuinely background $\aHo,\aH[(\Ny)]$ used later purely
to define the comoving coordinate $\yco$
(Section~\ref{sec:coordinate}), which necessarily stay
tied to a single global choice since they set up one shared comoving
labelling for the whole configuration -- see the discussion following
\eqref{eq:gradient-term-y}.
\end{panel}


Define the shell-averaged field for the set of comoving patches making up
shell $i$ as
\begin{equation}
    \field_i(\Ny) \equiv \frac{1}{\Nlay{i}} \sum_{\vect{x} \in i} \field(\vect{x}, \Ny) ,
\end{equation}
where $\Nlay{i}$ is the number of coarse-graining patches contained in
shell $i$. Averaging \eqref{eq:patch-langevin} over the shell, and
dropping the difference between $\langle \Vpot'(\field)\rangle$ and
$\Vpot'(\field_i)$ as a higher-order mean-field correction (suppressed by
$\Vpot''' \times \mathrm{Var}[\field]/\Nlay{i}$, itself already
$\Or(1/\Nlay{i})$ -- and, by the same argument, similarly approximating
$H^2,\epsSR$ evaluated on the shell by their value at the shell-mean
$\field_i,\mom_i$), gives
\begin{subequations}
\begin{align}
    \frac{\d \field_i}{\d \Ny}
    & =
    \mom_i + \overline{\noise{\field}}^{(i)} ,
    \\
    \frac{\d \mom_i}{\d \Ny}
    & =
    -\big(3-\epsSR(\field_i,\mom_i)\big)\mom_i - \frac{\Vpot'(\field_i)}{\Hub^2(\field_i,\mom_i)}
    + \frac{1}{\sfactor(\Ny)^2\Hub^2(\field_i,\mom_i)}\big\langle \grad^2\field \big\rangle_i
    + \overline{\noise{\mom}}^{(i)} .
\end{align}
\end{subequations}

\subsection{Central limit suppression of the shell-averaged noise}

The variance of the shell-averaged noise is
\begin{equation}
    \mathrm{Var}\big[\overline{\noise{\field}}^{(i)}\big]
    =
    \frac{1}{\Nlay{i}^2} \sum_{\vect{x},\vect{y} \in i}
    \Pspec{\field\field} \, \delta_D(\Ny-\Ny') \, \sinc(k_\sigma|\vect{x}-\vect{y}|) .
\end{equation}
If shell $i$ has linear comoving extent much larger than the
coarse-graining radius $k_\sigma^{-1}$, the double sum is dominated by
$\Or(\Nlay{i})$ correlated pairs out of $\Nlay{i}^2$ total, giving the
central-limit-theorem result
\begin{equation}
    \mathrm{Var}\big[\overline{\noise{\field}}^{(i)}\big]
    \; \approx \;
    \frac{c \, \Pspec{\field\field}}{\Nlay{i}} \, \delta_D(\Ny-\Ny') ,
    \label{eq:CLT-suppression}
\end{equation}
for some $\Or(1)$ geometric constant $c$, and similarly for
$\overline{\noise{\mom}}^{(i)}$.

\subsection{Inter-shell correlation is nearest-neighbour and subleading}

The covariance between the averaged noise in shells $i$ and $j$ is
\begin{equation}
    \mathrm{Cov}\big[\overline{\noise{\field}}^{(i)}, \overline{\noise{\field}}^{(j)}\big]
    =
    \frac{1}{\Nlay{i}\Nlay{j}} \sum_{\vect{x}\in i, \vect{y} \in j}
    \Pspec{\field\field} \, \delta_D(\Ny-\Ny') \, \sinc(k_\sigma|\vect{x}-\vect{y}|) .
\end{equation}
Since the sinc kernel connects only points within the coarse-graining
radius $k_\sigma^{-1}$, and the shell-to-shell separation is set by the
comoving Hubble radius --- parametrically larger than $k_\sigma^{-1}$
because $\sigma \ll 1$ --- this covariance vanishes for non-adjacent shells,
and for adjacent shells is suppressed both by the CLT factor
$1/\sqrt{\Nlay{i}\Nlay{j}}$ and by the fraction, $\Or(\Nlay{i}^{-1/3})$, of a
shell's patches that lie within the coarse-graining radius of its boundary.

\begin{panel}
Counting more carefully: the correlated part of the double sum has
$\Or(\Nlay{i}^{2/3})$ contributing pairs (those within a coarse-graining
length of the shared boundary) rather than $\Or(\Nlay{i}^2)$. Relative to
the leading CLT term \eqref{eq:CLT-suppression}, the correlation correction
within a single shell is therefore suppressed by an additional power
$\Nlay{i}^{-1/3}$. The inter-shell noise correlation is thus a genuinely
subleading effect once shells are averaged; the surviving inter-shell
coupling that matters at leading order is the \emph{deterministic} gradient
(Laplacian) term, discussed next.
\end{panel}

\subsection{The mean-field gradient (Laplacian) coupling}

The shell-averaged Laplacian, for a spherically decomposed field, is the
standard radial Laplacian,
\begin{equation}
    \big\langle \grad^2 \field \big\rangle_\mathrm{shell}
    \; \approx \;
    \frac{1}{\rad^2} \partial_\rad\!\big( \rad^2 \partial_\rad \field \big) .
    \label{eq:radial-laplacian-r}
\end{equation}
This deterministic coupling survives even in the strict large-shell limit
in which the noise correlation of the previous subsection vanishes: it is
the leading $\Or(\sigma^2)$ term in the gradient expansion and carries no
extra CLT suppression.

% ======================================================================
\section{Comoving shell geometry}
\label{sec:geometry}
% ======================================================================

Label shells by the step $i$ at which they detach from the active core.
With scale factor normalized so that $\sfactor(\Ninit) = \sfactor_0$, the
outer radius of shell $i$ is the comoving Hubble radius at step $i$,
\begin{equation}
    \rad_i = \frac{1}{\sfactor(i)\Hub(i)} ,
\end{equation}
and the shell width is $\Delta_i = \rad_i - \rad_{i+1}$. Because
$\d(\aH[^{-1}])/\d\Ny = -(1-\epsSR)\aH[^{-1}]$, the shell widths are not
uniform in comoving coordinates: outer shells (small $i$, born early, when
the comoving Hubble radius was larger) are wider than inner shells. The
comoving population of shell $i$ is fixed for all time once the shell is
created; what changes at later steps is only the field value it carries.

\begin{panel}
A discrete implementation along these lines --- a fixed number of
`active' stochastic shells (small enough that the CLT suppression
\eqref{eq:CLT-suppression} has not yet made the noise negligible), coupled
to an outer set of shells evolving under the deterministic gradient-coupled
equation only --- is a viable, and probably cheaper, alternative to the
continuum construction of Sections~\ref{sec:coordinate}--\ref{sec:collocation-scheme}.
The discrete and continuum treatments are two numerical schemes for the
\emph{same} underlying mean-field action, not different physical models. We
record only the continuum route in detail here, since it is the one
adopted for the first implementation, but the discrete route remains
available as a cross-check and may in the end prove more tractable; see
Section~\ref{sec:open-issues} and the companion planning document.
\end{panel}

% ======================================================================
% ======================================================================
\section{The onion coordinate $\yco$}
\label{sec:coordinate}
% ======================================================================

\subsection{Definition}

We wish to describe the geometry exterior to the horizon volume that
forms the core of the instanton. The inner edge of the domain is placed at
the core's own comoving horizon,
\begin{equation}
    \rH(\Ny) \equiv \frac{1}{a(\Ny)\,\Hub\big[\field(1,\Ny),\mom(1,\Ny)\big]} ,
    \label{eq:rH-definition}
\end{equation}
the local Friedmann-equation evaluation used throughout this construction,
now applied at the inner edge itself. This is a definition of the model,
not a derived fact: the ratio $\rho \equiv \rH(\Ny)/\rad(1,\Ny)$ of the
core's horizon to the radius of the inner edge is unity by construction.
The alternative of anchoring the inner edge to the background horizon
$\rHbg(\Ny) = 1/[a(\Ny)\Hub(\phinl,\pinl)]$ was considered and rejected;
see the panel on the anchor below.

The outer edge of the assembling overdensity is at a fixed comoving radius,
\begin{equation}
    \rout \equiv \frac{1+\alpha}{\aHo} ,
    \qquad \alpha \geq 0 ,
    \label{eq:rout-alpha}
\end{equation}
a small, controlled deformation of the naive choice $\rout=1/\aHo$ used in
an earlier draft, for reasons given in the panel below. We define the
radial variable $s \equiv \ln(\rad/\rout)$ --- motivated by the density
profile computed in the gradient-free approximation being close to a
power law in $\rad$, hence close to linear (and smooth) in $s$ --- and
write
\begin{equation}
    \Deltas(\Ny) \equiv \ln\!\left(\frac{\rout}{\rH(\Ny)}\right) ,
    \qquad
    \yco \equiv -\frac{2s}{\Deltas(\Ny)} - 1 ,
    \qquad
    \rad(\yco,\Ny) = \rout \exp\!\left[-\frac{(\yco+1)\Deltas(\Ny)}{2}\right] .
    \label{eq:y-definition}
\end{equation}
so that $\yco=-1$ is the fixed outer edge $\rout$ (which will turn out to
satisfy a Dirichlet boundary condition), and $\yco=+1$ is the current
horizon $\rH(\Ny)$ (which will turn out to satisfy a Neumann boundary
condition). The domain is therefore $\yco\in[-1,1]$ throughout, matching
the natural range of Legendre/Chebyshev collocation nodes
(Section~\ref{sec:collocation-basis}). Since $\rH(\Ninit)=1/\aHo$ exactly
(the core has not yet deviated from the noiseless trajectory at $\Ninit$),
\begin{equation}
    \Deltas(\Ninit) = \ln(1+\alpha) ,
    \label{eq:Deltas-init}
\end{equation}
strictly positive for $\alpha>0$: the resolved domain starts with a small
but finite log-radial extent, rather than as a single point.

Because $\rH$ is built from the core state, $\Deltas$ is an algebraic
function of $\field(1,\Ny)$, $\mom(1,\Ny)$ and $\Ny$,
\begin{equation}
    \Deltas(\Ny)
    =
    \ln(1+\alpha) + (\Ny-\Ninit)
    + \frac12\ln\frac{\Hub^2[\field(1,\Ny),\mom(1,\Ny)]}{\Hub^2[\phinl(\Ninit),\pinl(\Ninit)]} ,
    \label{eq:Deltas-closed}
\end{equation}
using $a(\Ny) = a(\Ninit)\e{\Ny-\Ninit}$. No differential equation is needed
to obtain it. Its rate of change along the solution is, however, fixed by
the chain rule and carries the gradient and noise forcing of the core; see
Section~\ref{sec:advection}.

\begin{panel}
\semibold{Why $\alpha>0$: the naive $\rout=1/\aHo$ choice is singular at $\Ninit$.}
With $\alpha=0$, $\Deltas(\Ninit)=0$ exactly, and the Laplacian coefficients
$4/\Deltas^2$, $2/\Deltas$ (\eqref{eq:Lop-definition}), the advection
coefficient $A\propto1/\Deltas$, and the noise dilution
$\ncount\propto\Deltas$ (Section~\ref{sec:noise-y}, so
$\Dnoise{ij}\propto1/\Deltas$) all diverge there. A Frobenius-type check of
the forward sector shows that the leading $1/\Deltas$ divergences in the
gradient and advection terms cancel against each other. Writing
$\field-\phinl=\Deltas(\Ny)\,h(\yco)+\Or(\Deltas^2)$ and
$\mom-\pinl=p_0(\yco)+\Or(\Deltas)$, and using that the gradient prefactor
$\gpref$ of \eqref{eq:gradient-term-y} equals one at $\Ninit$ up to
$\Or(\alpha)$, the $\Or(1/\Deltas)$ part of the $\mom$ equation is
\begin{equation}
    4\,h''(\yco) + (\yco+1)\big(1-\epscore\big)\,p_0'(\yco) = 0 ,
    \label{eq:frobenius-forward}
\end{equation}
a consistent leading order for the forward sector. The singularity is an
artefact of rescaling a shrinking physical domain onto the fixed range
$[-1,1]$: in $s$ the problem is regular on the shrinking interval
$[-\Deltas(\Ny),0]$. The response-field sector, which carries boundary data
from $\Nfinal$ and couples through the equally divergent
$\Dnoise{ij}\propto1/\Deltas$ noise terms, has not been analysed in the
same way. The measure convention of Section~\ref{sec:msr-action} removes
exactly one $1/\Deltas$ from the response equations (the one that came from
varying the measure) and restores parity between the two sectors in the
singular limit, which is the precondition for the response-sector indicial
analysis. It does not remove the geometric singularities: the advection
$1/\Deltas$, the Laplacian $1/\Deltas^2$ and the noise dilution $1/\Deltas$
are the same in every convention. The indicial analysis has not been done;
a back-of-envelope attempt did not obviously close.

Rather than complete that analysis, decouple $\rout$ from $\rH(\Ninit)$ by
the small amount $\alpha$ in \eqref{eq:rout-alpha}. Since
$\Deltas_\alpha(\Ny) = \ln(1+\alpha) + \Deltas_0(\Ny)$ (with $\Deltas_0$ the
original, unregularized quantity) is just an additive constant shift, every
divergent coefficient above is regularized identically, with no change to
any functional form --- and since the computed profiles already have
$\Deltas_0(\Nfinal)\sim9$--$21$, the shift is negligible almost immediately
after $\Ninit$ for any small $\alpha$, so late-time behaviour is
essentially unaffected. The physical content added is minimal: the thin
shell between $\rH(\Ninit)$ and the new $\rout$ is asserted to sit exactly
on the noiseless trajectory, extending an approximation already made
everywhere beyond $\rout$ slightly further in, not introducing a new kind
of assumption. Smaller $\alpha$ makes the (now finite) coefficients larger
near $\Ninit$, so there is a real stiffness trade-off in the choice; the
recommended approach is to run at a spread of $\alpha$ values (e.g.\
geometrically spaced) and check for the expected negligible sensitivity,
extrapolating to $\alpha\to0$ if any residual dependence remains, rather
than committing to a single value a priori.

Even if the indicial analysis closes, $\alpha$ is retained for two further
reasons. At $\Deltas(\Ninit)=0$ the resolved domain has zero log-radial
extent, so there is nothing for a boundary value problem in $\yco$ to
resolve. And near $\Deltas=0$ the individually divergent coefficients
suffer catastrophic cancellation in floating point even where their sum is
finite. $\alpha$ is therefore a deliberate numerical safeguard, not a
stand-in for an unverified cancellation.
\end{panel}

\begin{panel}
\semibold{This was the previously-rejected alternative --- now viable.} An
earlier version of this document considered exactly this construction
(anchoring the inner edge to the current horizon) and rejected it, on the
grounds that the resulting $\Ny$-dependent operator has no time-independent
eigenbasis, forcing re-diagonalization at every step. That objection is
specific to an \emph{explicit spectral method} (solving for coefficients
in an eigenfunction expansion); it does not apply to a \emph{collocation}
method, where the coupling introduced by an $\Ny$-dependent operator is no
worse in kind than the coupling already present from the potential and
noise terms (Section~\ref{sec:collocation-scheme}). Having settled on
collocation for other reasons, the original objection to this coordinate
no longer holds, and it directly resolves the sub-horizon-noise problem
(previously Section~13.1 in an earlier draft) by construction: there is no
region interior to the current horizon left in the domain to source
spurious independent fluctuations.
\end{panel}

\begin{panel}
\semibold{Core anchor versus background anchor.} The coordinate labels
$\rad$, $s$, $\yco$ are comoving labels of the unperturbed metric. On
14 July 2026 it was therefore proposed to anchor the inner edge to the
background horizon, $\Deltas = \ln(\rout/\rHbg)$, with the background
$\epsSR$ in the advection coefficient. That choice makes $\Deltas$ an
external function of $\Ny$, removes every state dependence from the
scaffold, and drops one term from the response-sector ledger. It was
rejected on 16 July for the following reason. With a background anchor
the core's own horizon sits at $\rho = \rH/\rHbg = \Hub_{\rm bg}/\Hub_{\core} < 1$
for an overdensity, and the characteristic analysis of
Section~\ref{sec:bcs} gives core speeds $(2/\Deltas)[-(1-\epsSR^{\rm nl}) \pm \rho]$.
The number of incoming characteristics at the core is then two for
$\epsSR^{\rm nl} < 1-\rho$, one for $1-\rho < \epsSR^{\rm nl} < 1+\rho$, and zero
above. Two small quantities compete ($\epsSR \sim 0.009$ against
$\delta\Hub/\Hub \sim 0.002$ at $\dNstar = 0.2$ for a quadratic potential),
the margin shrinks with $\dNstar$, and a zero crossing during a solve would
change the number of admissible boundary conditions. The physical reading
is that the domain is then too small: shells just inside the background
horizon are already outside their own (smaller) horizon and belong in the
domain, and their entry is the second incoming characteristic.

With the core anchor $\rho\equiv1$ by construction and the window is
$0<\epscore<2$ unconditionally. Two checks support it. In exact de Sitter
an ingoing null ray launched from the horizon stays on the horizon, so its
rate relative to the inner edge must vanish; the outgoing characteristic
speed $2\epscore/\Deltas$ does so. The incoming rate tends to $4/\Deltas$,
a relative rate of $2$ in $\ln\rad$, as it should for the outgoing ray in
de Sitter. The background anchor passes neither check. The price is that
$\Deltas$, $\advcoef$, the measure and $\ncount$ are all functions of the
core state, which makes the response sector long to derive by hand
(Section~\ref{sec:instanton-eqs}). The decision was that this is a
self-consistent choice that is neither right nor wrong: it is the model.
Whether it is adequate could only be tested against lattice simulations.
\end{panel}

\subsection{The advection term}
\label{sec:advection}

The map \eqref{eq:y-definition} is time-dependent at fixed $\rad$, since
$\Deltas(\Ny)$ evolves. Writing $s=-(\yco+1)\Deltas(\Ny)/2$ at fixed
$\rad$ (i.e.\ fixed $s$) and differentiating,
\begin{equation}
    \left.\frac{\partial \yco}{\partial \Ny}\right|_\rad
    =
    -\frac{(\yco+1)}{\Deltas(\Ny)}\,\dot{\Deltas}(\Ny) ,
    \qquad
    \dot{\Deltas} \equiv \frac{\d\Deltas}{\d\Ny}
    = 1 + \frac{\d}{\d\Ny}\ln\Hub\big[\field(1,\Ny),\mom(1,\Ny)\big] ,
    \label{eq:Deltas-derivative}
\end{equation}
where the second equality follows from \eqref{eq:Deltas-closed}. Define the
advection coefficient
\begin{equation}
    \advcoef(\yco,\Ny) \equiv \frac{\yco+1}{\Deltas(\Ny)}\,\dot{\Deltas}(\Ny) ,
    \label{eq:advcoef-def}
\end{equation}
so that
\begin{equation}
    \left.\frac{\partial}{\partial \Ny}\right|_\rad
    =
    \left.\frac{\partial}{\partial \Ny}\right|_\yco
    -
    \advcoef(\yco,\Ny)\,\partial_\yco .
    \label{eq:advection-operator}
\end{equation}
$\advcoef$ vanishes at $\yco=-1$ (the fixed outer edge is correctly
undisturbed) and is largest at $\yco=+1$ (the moving horizon).

\para{The time derivative of $\Deltas$ is derived, not posited}
Because $\Deltas$ is an algebraic function of the core state, its rate of
change is fixed by the chain rule. With $\Hub^2=\Vpot/(3-\epsSR)$ and
$\epsSR=\mom^2/2$ in reduced Planck units,
\begin{equation}
    \dot{\Deltas}
    =
    1 + \frac12\left[\frac{\Vpot'(\field_\core)}{\Vpot(\field_\core)}\,\dot\field_\core
    + \frac{\mom_\core\,\dot\mom_\core}{3-\epscore}\right] ,
    \label{eq:Deltas-dot-chain}
\end{equation}
where the subscript $\core$ denotes evaluation at $\yco=+1$. The rates
$\dot\field_\core$, $\dot\mom_\core$ are the core rows of the instanton
equations of Section~\ref{sec:instanton-eqs}, which are not the
homogeneous equations. Write those rows as
\begin{subequations}
\label{eq:core-rows}
\begin{align}
    \dot\field_\core
    & = \mom_\core + \advcoef_\core\,(\partial_\yco\field)_\core + \Fcore^{\field} ,
    \\
    \dot\mom_\core
    & = -(3-\epscore)\mom_\core - \frac{\Vpot'(\field_\core)}{\Hub^2_\core}
    + \advcoef_\core\,(\partial_\yco\mom)_\core + \Fcore^{\mom} ,
\end{align}
\end{subequations}
with $\advcoef_\core \equiv \advcoef(1,\Ny) = 2\dot{\Deltas}/\Deltas$, and
$\Fcore^{\field}$, $\Fcore^{\mom}$ the non-advective forcing at the core:
the noise sourcing in both rows, and the gradient term in the $\mom$ row.
Substituting \eqref{eq:core-rows} into \eqref{eq:Deltas-dot-chain} and
collecting the terms proportional to $\dot{\Deltas}$ gives a linear
equation for $\dot{\Deltas}$, with solution
\begin{equation}
    \dot{\Deltas}
    =
    \frac{1 - \epscore
    + \dfrac{\Vpot'(\field_\core)}{2\Vpot(\field_\core)}\,\Fcore^{\field}
    + \dfrac{\mom_\core\,\Fcore^{\mom}}{2(3-\epscore)}}{1-\Kcoef} ,
    \qquad
    \Kcoef \equiv \frac{1}{\Deltas}\left[
    \frac{\Vpot'(\field_\core)}{\Vpot(\field_\core)}(\partial_\yco\field)_\core
    + \frac{\mom_\core\,(\partial_\yco\mom)_\core}{3-\epscore}
    \right] .
    \label{eq:Deltas-dot-closed}
\end{equation}
This is a closed form, evaluated from the state at each instant; no extra
state variable and no iteration are needed. Under the Neumann condition
$(\partial_\yco\field)_\core=0$ of Section~\ref{sec:bcs} the first term of
$\Kcoef$ vanishes identically.

\begin{calculation}
\semibold{Collapse to $1-\epscore$ on the homogeneous equations.}
Substitute the homogeneous rates $\dot\field=\mom$,
$\dot\mom=-(3-\epsSR)\mom-\Vpot'/\Hub^2$ into \eqref{eq:Deltas-dot-chain}:
\begin{equation}
    \frac12\left[\frac{\Vpot'}{\Vpot}\mom - \mom^2 - \frac{\mom\Vpot'}{(3-\epsSR)\Hub^2}\right]
    = \frac12\left[\frac{\Vpot'}{\Vpot}\mom - \mom^2 - \frac{\Vpot'}{\Vpot}\mom\right]
    = -\frac{\mom^2}{2} = -\epsSR ,
\end{equation}
using $(3-\epsSR)\Hub^2=\Vpot$. Hence $\dot{\Deltas}=1-\epscore$, the
familiar identity $\d\ln(a\Hub)/\d\Ny=1-\epsSR$. In
\eqref{eq:Deltas-dot-closed} the same cancellation produces the leading
$1-\epscore$; the remaining terms are the advective and non-advective
forcing of the core row.
\end{calculation}

\begin{panel}
\semibold{Why the collapse is not available here.} The identity
$\epsSR=\mom^2/2$ is exact in a homogeneous FRW patch, and
$\dot{\Deltas}=1-\epscore$ looks like a constitutive statement. It is not.
Its proof consumes the homogeneous $\dot\mom$ equation, and the core row of
the onion model is not homogeneous: it carries the gradient term, the
advection term and the noise sourcing. An earlier version of these notes
paired the algebraic $\Deltas$ of \eqref{eq:Deltas-closed} with the
background rate $\dot{\Deltas}=1-\epscore$. That pairing is inconsistent:
differentiating the closed form does not give $1-\epscore$ off the
background, and the mismatch is precisely a drift of the inner edge away
from the horizon it is meant to track. The failure is seductive because the
identity is exact rather than slow-roll; it is nevertheless a consequence
of an equation of motion, not a definition.

Three remarks on \eqref{eq:Deltas-dot-closed}. First, $\Kcoef\to0$ at a
converged, regular solution, where the $\yco$-derivatives at the core
vanish; mid-iteration it need not. Second, the numerator correction does
not vanish at convergence: $\Fcore^{\mom}$ contains the Laplacian at the
core, which is maximal at a density peak, and the noise sourcing, which is
nonzero on the instanton. The core's Hubble rate genuinely responds to the
gradient, and a grid anchored to the core's horizon must follow it. Third,
$\Kcoef=1$ is a genuine singularity of the coordinate map, where the
horizon rate becomes infinitely sensitive to the profile; it should be
guarded against rather than discovered.
\end{panel}

\begin{todo}
Whether a numerical boundary penalty (the SAT of
Section~\ref{sec:collocation-basis}) belongs in $\Fcore^{\mom}$ when
$\dot{\Deltas}$ is evaluated is open. The argument against is that the
penalty is a stabilizer rather than physics, vanishes at convergence, and
keeps the closed form free of a second linear loop through the penalty
strength. See \texttt{RECONSTRUCTION.md}, Open~3.
\end{todo}

\subsection{The gradient operator in $\yco$}
\label{sec:gradient-operator}

Using the standard log-radial identity $\nabla^2\field = \rad^{-2}(\field_{ss}+\field_s)$
and $\partial_s = -2 [ \Deltas(\Ny) ]^{-1} \partial_\yco$ (constant at fixed
$\Ny$, so no further chain-rule correction is needed for the second
derivative),
\begin{equation}
    \nabla^2\field
    =
    \rout^{-2}\,\e{\Deltas(\Ny)}\,\e{\Deltas(\Ny)\,\yco}
    \left[\frac{4}{\Deltas(\Ny)^2}\partial_\yco^2\field - \frac{2}{\Deltas(\Ny)}\partial_\yco\field\right]
    \equiv
    \rout^{-2}\,\Lop\field ,
    \label{eq:Lop-definition}
\end{equation}
using $1/\rad(\yco,\Ny)^2 = \rout^{-2}\,\e{\Deltas(\Ny)}\e{\Deltas(\Ny)\yco}$
from \eqref{eq:y-definition} --- note this uses $\rout$ itself, \emph{not}
$\aHo$: with $\alpha>0$ (\eqref{eq:rout-alpha}) these are no longer reciprocals of each
other, $\rout=(1+\alpha)/\aHo$, and using $\aHo^2$ here (as an earlier
draft did, implicitly assuming $\alpha=0$) would be wrong by a factor of
$(1+\alpha)^2$. This operator
carries an explicit first-derivative (drift) term and depends on $\Ny$
through $\Deltas(\Ny)$, not just through an overall $\yco$-independent
prefactor. The gradient term in the $\mom$ equation of motion is
\begin{equation}
    \frac{1}{\aHloc^2}\big\langle\grad^2\field\big\rangle
    =
    \frac{1}{\rout^2\aHloc^2}\,\Lop\field ,
    \label{eq:gradient-term-y}
\end{equation}
with $\aHloc$ the same local quantity as before (Section~\ref{sec:noise-y}).

\begin{panel}
\semibold{The prefactor as a ``local $\Deltas$''.} Since
$\Deltas(\Ny)=\ln(\rout/\rH(\Ny))=\ln\big(\rout\,a(\Ny)\,\Hub[\field(1,\Ny),\mom(1,\Ny)]\big)$
(core $\Hub$, Section~\ref{sec:coordinate}), define the local analogue
\begin{equation}
    \Deltas_{\rm loc}(\yco,\Ny)
    \equiv
    \ln\big(\rout\,a(\Ny)\,\Hub[\field(\yco,\Ny),\mom(\yco,\Ny)]\big)
    =
    \ln\big(\rout\,\aHloc(\yco,\Ny)\big) ,
\end{equation}
same functional form as $\Deltas(\Ny)$ with the core Hubble rate replaced
by the local (shell-$\yco$) one. Then directly
$\rout^{-2}\aHloc^{-2}=\e{-2\Deltas_{\rm loc}(\yco,\Ny)}$, so the prefactor
in \eqref{eq:gradient-term-y} is exactly a local-$\Deltas$ exponential.
Consistency check: at $\yco=1$, $\aHloc(1,\Ny)=1/\rH(\Ny)$ exactly
(Section~\ref{sec:noise-y}), so $\Deltas_{\rm loc}(1,\Ny)=\Deltas(\Ny)$,
recovering the global definition at the core as it must. This is purely a
notational simplification for the \emph{prefactor} multiplying $\Lop$; it
must \emph{not} be confused with the (genuinely global, core-based)
$\Deltas(\Ny)$ baked into $\Lop$ itself via \eqref{eq:Lop-definition},
which comes from the $\rad\to\yco$ coordinate map and is necessarily
shared by every shell, since $\yco$ is a single coordinate label common to
the whole domain.
\end{panel}

\subsubsection{Self-adjoint measure}

Writing $\Lop$ in $p(\yco)\field''+q(\yco)\field'$ form (dropping the
$\yco$-independent prefactor $\rout^{-2}\e{\Deltas}\cdot4/\Deltas^2$ for
clarity, and absorbing it back at the end), $p=\e{\Deltas\yco}$,
$q=-\tfrac12\Deltas\,\e{\Deltas\yco}$. The Sturm--Liouville reduction
$\measure^{-1}(\measure p\,\field')'$ requires $\measure'/\measure=(q-p')/p$;
here $(q-p')/p = -\tfrac32\Deltas(\Ny)$, a genuine constant (not a function
of $\yco$), so the reduction is exact and closed-form:
\begin{equation}
    \measure(\yco,\Ny) = \e{-\frac32\Deltas(\Ny)\,\yco} ,
    \qquad
    \Lop\field
    \;\propto\;
    \e{\Deltas(\Ny)\yco}\left(\e{-\frac32\Deltas(\Ny)\yco}\right)^{-1}\frac{\d}{\d\yco}\!\left(\e{-\frac12\Deltas(\Ny)\yco}\,\frac{\d\field}{\d\yco}\right) .
    \label{eq:self-adjoint-measure}
\end{equation}
As a consistency check, $\rad(\yco,\Ny)^3 \propto \e{-\frac32\Deltas(\Ny)\yco}$
--- $\measure(\yco,\Ny)$ is, up to normalization, exactly the physical
comoving volume measure, as it must be: $\nabla^2$ is self-adjoint with
respect to $\rad^2\,\d\rad$ in \emph{any} radial parametrization, by
Green's theorem, so recovering it here from the Sturm--Liouville reduction
is a check on the algebra rather than a new input.

\begin{panel}
\semibold{Flat-measure (Schrödinger normal) form: explored and set aside.}
The further Liouville substitution that removes $q(\yco)$ \emph{and}
brings the measure to flat $\d\yco$ (rather than the weighted
$\measure(\yco,\Ny)\,\d\yco$ above) was worked through explicitly. It
succeeds algebraically --- the residual potential term vanishes exactly,
an exactly-solvable special case, since the exponential coefficients here
play the same role as in an Euler-type ODE --- but the new independent
variable $x(\yco)=\tfrac{2}{\Deltas(\Ny)}(1-\e{-\Deltas(\Ny)\yco/2})$ maps
$\yco\in[-1,1]$ onto a badly asymmetric interval whose width grows like
$\e{\Deltas(\Ny)/2}$: for the wide ($\dNstar=13$) profile already computed
in the gradient-free approximation, $\Deltas\approx20.7$ gives a domain
width $\sim3\times10^3$; even the narrower ($\dNstar\approx1.5$) case gives
$\Deltas\approx9.2$ and a width $\sim22$. This reintroduces, in the
transformed coordinate, essentially the same exponential compression the
logarithmic radial variable was adopted to remove. \semibold{Conclusion:}
retain the \emph{weighted} measure \eqref{eq:self-adjoint-measure} and
build the self-adjoint discrete structure (Section~\ref{sec:collocation-scheme})
directly for it, rather than transforming to a flat measure.
\end{panel}

This removes the previous $\S4.4$ exponential-mode-resolution concern
structurally, not by mitigation: the logarithmic radial variable is
precisely what turns the exponentially-compressed near-horizon structure
(in linear $\rad$) into $\Or(1)$ structure in $\yco$ --- consistent with
the density profile already computed in the gradient-free approximation
being close to a power law in $\rad$, hence smooth and close to linear in
$s=\ln(\rad/\rout)$. A first implementation should confirm this
empirically once collocation-point convergence can be checked, but there
is no longer a structural reason to expect the mode/point count needed to
grow exponentially with elapsed $e$-folds.

% ======================================================================
\section{Noise statistics in $\yco$ coordinates}
\label{sec:noise-y}
% ======================================================================


The number of coarse-graining (current-Hubble-volume) patches per unit
$\yco$-shell can be computed purely in comoving terms. From
\eqref{eq:y-definition}, $\d\rad/\d\yco = -\tfrac12\Deltas(\Ny)\,\rad(\yco,\Ny)$
at fixed $\Ny$, so a shell of coordinate width $\d\yco$ has comoving volume
$4\pi\rad^2|\d\rad| = 2\pi\Deltas(\Ny)\,\rad(\yco,\Ny)^3\,\d\yco$. The
comoving volume of the current Hubble patch \emph{local to that shell} ---
which is what matters, since noise is a spatially local process --- is
$\tfrac43\pi/\aHloc^3$, using the same local $\aHloc$ as the gradient term
(Section~\ref{sec:gradient-operator}). The number of patches per unit
$\yco$ is the ratio of these volumes,
\begin{equation}
    \ncount(\yco,\Ny)
    =
    \frac32\Deltas(\Ny)\,\big[\rad(\yco,\Ny)\,\aHloc\big]^3
    =
    \frac32\Deltas(\Ny)\,\big[\rout\,\aHloc\big]^3 \, \e{-\frac32(\yco+1)\Deltas(\Ny)} ,
    \label{eq:ncount}
\end{equation}
using $\rad(\yco,\Ny) = \rout\,\e{-(\yco+1)\Deltas(\Ny)/2}$ directly from
\eqref{eq:y-definition} --- note this is written in terms of $\rout$, not
$\aHo$: since $\rout=(1+\alpha)/\aHo$, these differ once $\alpha>0$, and
the combination $\rout\aHloc$ is the one that appears without extra
factors of $(1+\alpha)$. As before, the $\aH$ inside $\aHloc$ is the local
quantity used throughout Section~\ref{sec:instanton-eqs}, while $\rout$
retains its fixed, coordinate-defining role. At the outer edge
($\yco=-1$, where $\field=\phinl$ is pinned), $\aHloc\to a(\Ny)\Hub[\phinl(\Ny),\pinl(\Ny)]$,
generically $\gg\aHo$ (hence $\rout\aHloc\gg1$, since $\rout=\Or(1/\aHo)$)
at late $\Ny$, so $\ncount$ is large there (many patches, as expected). At
the horizon itself ($\yco=+1$), $\aHloc$ is by \eqref{eq:rH-definition}
exactly $1/\rH(\Ny)$, so $\rout\aHloc = \rout/\rH(\Ny) = \e{\Deltas(\Ny)}$
directly from the definition of $\Deltas(\Ny)$, and the two exponential
factors cancel: $\ncount(1,\Ny) = \tfrac32\Deltas(\Ny)$
--- an $\Or(10)$, not $\Or(1)$, quantity for the profiles already computed
($\Deltas\sim9$--$21$ depending on $\dNstar$): the exterior-only domain no
longer approaches the single-patch limit $\ncount\to1$ at its inner edge,
because the excised sub-horizon interior --- not the last resolved shell
--- is where that single remaining patch now lives, outside the domain
entirely.

\begin{panel}
\semibold{$\ncount(\yco,\Ny)$ as a density, not a shell count.}
\eqref{eq:ncount} is a \emph{density} $\d N_{\mathrm{patches}}/\d\yco$, not
the (dimensionless) patch count in a shell of finite width: the count in a
shell of coordinate width $\d\yco$ is $\ncount(\yco,\Ny)\,\d\yco$, and it
is the density that enters $\Dnoise{\field}=\Pspec{\field\field}/\ncount$
below, consistently with the $\delta(\yco-\yco')$ noise correlation of
\eqref{eq:langevin-y}, which requires a density on the right-hand side
rather than a quantity that vanishes as $\d\yco\to0$.
\end{panel}

\begin{panel}
\semibold{Geometric dilution versus local amplitude.} $\ncount(\yco,\Ny)$
above is \emph{not} purely a coordinate/gauge quantity, unlike $\aHo,\rH(\Ny)$
as used to define $\yco$ itself: it uses $\aHloc$, the same local $aH$ as
the gradient term, because ``how big is the current Hubble patch at this
shell'' is a genuinely local physical statement, not a bookkeeping choice
about the comoving coordinate map. The single-patch \emph{amplitude}
entering the average, $\Dcm{ij}(\field,\mom)$, is a separate object with
its own local dependence: the codebase already treats the diffusion matrix
as a local function of the phase-space point --
\texttt{DiffusionModel.D\_matrix(phi,pi)} returns $(\Dcm{11},\Dcm{12},\Dcm{22})$
evaluated at $(\field,\mom)$, exactly as \texttt{AbstractPotential} does for
$V$, $H^2$, $\epsilon_1$ -- and the existing single-trajectory equations of
motion (\texttt{ComputeTargets/FullInstanton.py}) use it that way,
e.g.\ $\dot\field = \mom + 2\Dcm{11}\rfield + 2\Dcm{12}\rmom$.
\end{panel}

Using $\Pspec{fg}\atloc = 2\,\Dcm{fg}\atloc$ (the identification implicit in
$\Pspec{}/2 = H^2/8\pi^2$ in the slow-roll, decoupled-diffusion case), the
correctly local shell-averaged noise variance/covariance densities are
\begin{equation}
    \Dnoise{\field}(\yco,\Ny) \equiv \frac{2\,\Dcm{11}\big(\field(\yco,\Ny),\mom(\yco,\Ny)\big)}{\ncount(\yco,\Ny)} ,
    \qquad
    \Dnoise{\mom}(\yco,\Ny) \equiv \frac{2\,\Dcm{22}\big(\field(\yco,\Ny),\mom(\yco,\Ny)\big)}{\ncount(\yco,\Ny)} ,
    \label{eq:Dnoise-diag}
\end{equation}
together with the cross term,
\begin{equation}
    \Dnoise{\field\mom}(\yco,\Ny) \equiv \frac{2\,\Dcm{12}\big(\field(\yco,\Ny),\mom(\yco,\Ny)\big)}{\ncount(\yco,\Ny)} .
    \label{eq:Dnoise-cross}
\end{equation}
All three are evaluated pointwise on the same $\yco$-grid, using the same
reconstructed $\field(\yco,\Ny)$, $\mom(\yco,\Ny)$, by the same procedure
used for every other nonlinear, solution-dependent coefficient in this
scheme (Section~\ref{sec:collocation-scheme}). Since $\ncount$ no longer
approaches $1$ anywhere in the resolved domain, treating
$\Dnoise{\field}(\yco,\Ny)$ as a pointwise, $\delta(\yco-\yco')$-correlated
density is on safer footing throughout $\yco\in[-1,1]$ than it was for the
previous, interior-covering coordinate.
% ======================================================================

\subsection{The measure selected by the noise}
\label{sec:noise-measure}

The covariance \eqref{eq:Dnoise-diag} was obtained by counting patches.
The same result follows from an explicit angle average, and that
derivation shows with respect to which volume measure the noise is white.
This fixes a normalization that the self-adjointness argument of
Section~\ref{sec:gradient-operator} leaves free, and it determines the
measure to be used in the MSR action (Section~\ref{sec:msr-action}).

\begin{calculation}
\semibold{Angle-averaged noise correlator.} At separations larger than the
coarse-graining length the sinc kernel of \eqref{eq:noise-covariance} acts
as a delta function, so take the patch-level noise to be spatially white,
\begin{equation}
    \big\langle\noise{\field}(\vect{x},\Ny)\,\noise{\field}(\vect{x}',\Ny')\big\rangle
    = \mathcal{A}\,\delta^{(3)}(\vect{x}-\vect{x}')\,\delta_D(\Ny-\Ny') ,
\end{equation}
with $\mathcal{A}$ fixed by requiring that the average over one Hubble
patch, of comoving volume $\vH = \tfrac43\pi/\aHloc^3$, reproduce the
single-patch variance: $\langle\bar\xi_{\rm patch}^2\rangle = \mathcal{A}/\vH = 2\Dcm{11}$,
so $\mathcal{A} = 2\Dcm{11}\,\vH$. The Langevin variable of the onion model
is the angle average on the shell, $\bar\xi(\rad) = (4\pi)^{-1}\int\d\Omega\,\noise{\field}$.
Using $\delta^{(3)}(\vect{x}-\vect{x}') = \rad^{-2}\delta(\rad-\rad')\,\delta^{(2)}(\hat{\vect{n}}-\hat{\vect{n}}')$,
\begin{equation}
    \big\langle\bar\xi(\rad)\,\bar\xi(\rad')\big\rangle
    = \frac{\mathcal{A}}{(4\pi)^2}\int\d\Omega\,\d\Omega'\,
      \frac{\delta(\rad-\rad')}{\rad^2}\,\delta^{(2)}(\hat{\vect{n}}-\hat{\vect{n}}')
    = \frac{\mathcal{A}}{4\pi\rad^2}\,\delta(\rad-\rad')\,\delta_D(\Ny-\Ny') .
\end{equation}
Converting to $\yco$ with $|\d\rad/\d\yco| = \tfrac12\Deltas\,\rad$,
\begin{equation}
    \big\langle\bar\xi(\yco)\,\bar\xi(\yco')\big\rangle
    = \frac{\mathcal{A}}{2\pi\Deltas\,\rad^3}\,\delta(\yco-\yco')\,\delta_D(\Ny-\Ny')
    = \frac{\mathcal{A}}{\volel(\yco,\Ny)}\,\delta(\yco-\yco')\,\delta_D(\Ny-\Ny') ,
    \label{eq:noise-angle-averaged}
\end{equation}
where
\begin{equation}
    \volel(\yco,\Ny) \equiv \frac{\d V}{\d\yco} = 4\pi\rad^2\left|\frac{\d\rad}{\d\yco}\right|
    = 2\pi\rout^3\,\Deltas(\Ny)\,\e{-\frac32(\yco+1)\Deltas(\Ny)}
    \label{eq:volel-definition}
\end{equation}
is the physical comoving volume element per unit $\yco$. Finally
\begin{equation}
    \frac{\mathcal{A}}{\volel}
    = \frac{2\Dcm{11}\,\tfrac43\pi/\aHloc^3}{2\pi\Deltas\,\rad^3}
    = \frac{2\Dcm{11}}{\tfrac32\Deltas\,(\rad\,\aHloc)^3}
    = \frac{2\Dcm{11}}{\ncount(\yco,\Ny)} ,
\end{equation}
which is \eqref{eq:Dnoise-diag}. The amplitude quoted there is therefore
exactly right, with a flat $\delta(\yco-\yco')$.
\end{calculation}

The invariant object in \eqref{eq:noise-angle-averaged} is
$\delta(\yco-\yco')/\volel(\yco,\Ny)$, the delta function normalized
against $\volel\,\d\yco$: $\int\volel(\yco')\,\d\yco'\,\delta(\yco-\yco')/\volel(\yco') = 1$.
The noise is white with respect to the physical volume element, not with
respect to the self-adjoint measure $\measure=\e{-\frac32\Deltas\yco}$ of
\eqref{eq:self-adjoint-measure}. The two differ by the $\yco$-independent
factor
\begin{equation}
    \volel(\yco,\Ny) = \Cpref(\Ny)\,\measure(\yco,\Ny) ,
    \qquad
    \Cpref(\Ny) = 2\pi\rout^3\,\Deltas(\Ny)\,\e{-\frac32\Deltas(\Ny)} ,
    \label{eq:Cpref-definition}
\end{equation}
which self-adjointness cannot see (it depends only on the $\yco$ structure)
but which depends on $\Ny$ and so enters the MSR action through the time
derivative of the measure. The noise selects $\volel$.

\begin{panel}
\semibold{Noise density versus covariance, and two distinct volume factors.}
In \eqref{eq:noise-angle-averaged} the noise $\bar\xi(\yco)$ is a
density: the physical fluctuation received by a shell of coordinate width
$\d\yco$ is $\bar\xi\,\volel\,\d\yco$. Its covariance is
$\mathcal{A}\,\delta(\yco-\yco')/\volel$, a distribution; it is not the
smooth reciprocal of a precision. Contracting one factor of $\volel\,\d\yco$
against the delta gives
$\mathrm{Var}(\bar\xi\,\volel\,\d\yco) = \mathcal{A}\,\volel\,\d\yco$,
first order in the shell volume: the density covariance shrinks as
$1/\volel$ and the extensive per-shell covariance grows as $\volel$, both
from one delta on top of an intensive amplitude $\mathcal{A}$. When the
action is written with a precision kernel instead, $\mathcal{A}$ appears
inverted; the two descriptions are reciprocal, not different.

Two volume factors appear in the construction and must not be cancelled
against each other. The dilution $1/\ncount$ in \eqref{eq:Dnoise-diag}
converts per-patch noise to per-shell-mean noise and is a property of the
Langevin equation, independent of any action. The measure $\volel\,\d\yco$
is the quadrature weight of the action and also fixes the energy norm of
Section~\ref{sec:collocation-basis}. Both carry $\e{3\Deltas}$-type
factors and do different jobs. The statement that the Gaussian shift
$\noise{}-\im P/\mathcal{A}\to\noise{}$ leaves the path-integral measure
invariant is a leading-order statement: for a field-dependent amplitude the
Jacobian of that shift is the one-loop fluctuation determinant.
\end{panel}
\section{The Langevin system in $(\yco,\Ny)$}
\label{sec:langevin-y}
% ======================================================================

Collecting the results of Sections~\ref{sec:coordinate}--\ref{sec:noise-y},
the mean-field Langevin system is
\begin{subequations}
\label{eq:langevin-y}
\begin{align}
    \frac{\partial \field}{\partial \Ny}
    & =
    \mom + \advcoef(\yco,\Ny)\,\partial_\yco\field + \noise{\field}(\yco,\Ny) ,
    \\
    \frac{\partial \mom}{\partial \Ny}
    & =
    -\big(3-\epsSR\atloc\big)\mom - \frac{\Vpot'(\field)}{\Hub^2\atloc}
    + \frac{1}{\rout^2\aHloc^2} \Lop \field
    + \advcoef(\yco,\Ny)\,\partial_\yco\mom
    + \noise{\mom}(\yco,\Ny) ,
\end{align}
\end{subequations}
where the derivatives on the left are now at fixed $\yco$
(\eqref{eq:advection-operator} converts the original fixed-$\rad$ equations),
and the advection term $\advcoef\,\partial_\yco(\cdot)$ --- absent for the
previous coordinate, present here because the map itself is
$\Ny$-dependent (Section~\ref{sec:coordinate}) --- appears in \emph{both}
equations, since both were originally written at fixed $\rad$. As before,
$\Hub^2\atloc$ and $\epsSR\atloc$ denote $H_{\mathrm{sq}}[\field(\yco,\Ny),\mom(\yco,\Ny)]$
and $\epsilon[\field(\yco,\Ny),\mom(\yco,\Ny)]$, the same local
Friedmann-equation functions used elsewhere, distinct from $\aHo$ and the
core-anchored $\rH(\Ny)$ that define $\yco$ itself and $\ncount(\yco,\Ny)$. On the fixed rectangular domain $\yco\in[-1,1]$,
$\Ny\in[\Ninit,\Nfinal]$, the noises have covariance
\begin{equation}
    \big\langle \noise{\field}(\yco,\Ny)\noise{\field}(\yco',\Ny')\big\rangle
    =
    \Dnoise{\field}(\yco,\Ny)\,\delta_D(\Ny-\Ny')\,\delta(\yco-\yco') ,
\end{equation}
similarly for $\big\langle\noise{\mom}\noise{\mom}\big\rangle$ with
$\Dnoise{\mom}$, and
\begin{equation}
    \big\langle \noise{\field}(\yco,\Ny)\noise{\mom}(\yco',\Ny')\big\rangle
    =
    \Dnoise{\field\mom}(\yco,\Ny)\,\delta_D(\Ny-\Ny')\,\delta(\yco-\yco') ,
    \label{eq:noise-corr}
\end{equation}
where the $\delta(\yco-\yco')$ reflects the result of
Section~\ref{sec:shell-averaging} that inter-shell noise correlations are
subleading once the CLT suppression has been accounted for. Unlike the
previous coordinate, there is no longer any sub-horizon region within the
domain to worry about: $\yco=+1$ \emph{is} the current horizon by
construction, so the domain boundary and the physical boundary between
``many independent patches'' and ``at most one patch's worth of freedom''
coincide exactly, at every $\Ny$.

% ======================================================================
% ======================================================================
\section{The MSR action}
\label{sec:msr-action}
% ======================================================================

\subsection{Choice of measure}
\label{sec:measure-choice}

The MSR action integrates over $\yco$ against a reference density
$\varrho(\yco,\Ny)\,\d\yco$, and the response fields are densities
conjugate to $\field$, $\mom$ with respect to that measure. A choice of
$\varrho$ is a convention: densities related by
$\varrho\to\lambda(\yco,\Ny)\varrho$ give the same physics once the
response fields are rescaled by $1/\lambda$ and the noise amplitude by
$\lambda$. Three choices are natural,
\begin{equation}
    \varrho = 1 \ \text{(flat)} , \quad
    \varrho = \measure = \e{-\frac32\Deltas\yco} \ \text{(self-adjoint)} , \quad
    \varrho = \volel = \Cpref(\Ny)\,\measure \ \text{(volume element)} ,
\end{equation}
with $\Cpref$ as in \eqref{eq:Cpref-definition}. Self-adjointness of
$\Lop$ requires only $\varrho'/\varrho = -\tfrac32\Deltas$ at fixed $\Ny$,
so it selects $\measure$ up to the factor $\Cpref(\Ny)$ and cannot
distinguish the last two. The noise does: by Section~\ref{sec:noise-measure}
it is white with respect to $\volel\,\d\yco$. We write the action with
$\varrho=\volel$.

The three conventions differ in the response sector only through the
multiplicative coefficient of the response fields generated by the
measure, which Section~\ref{sec:instanton-eqs} shows to be
\begin{equation}
    \mathcal{Z}[\varrho] \equiv -\partial_\Ny\ln\varrho + \partial_\yco\advcoef + \advcoef\,\partial_\yco\ln\varrho ,
    \qquad
    \mathcal{Z}[1] = \frac{\dot{\Deltas}}{\Deltas} , \quad
    \mathcal{Z}[\measure] = \dot{\Deltas}\left[\frac{1}{\Deltas}-\frac32\right] , \quad
    \mathcal{Z}[\volel] = 0 ,
    \label{eq:Z-conventions}
\end{equation}
and in the form of the gradient operator acting on the response fields
($\Lop$ in the two weighted conventions, its flat-measure adjoint in the
flat one). The forward diffusion coefficient is $\Dnoise{ij}=2\Dcm{ij}/\ncount$
in the flat convention and $\varrho\,\Dnoise{ij}$ in the others. All three
share the same forward solution and the same on-shell action.

\begin{panel}
\semibold{The earlier action was a splice.} The 6 July version of
\eqref{eq:msr-action} carried the measure $\measure$ on the linear terms
and the flat-convention coefficient $\Dnoise{ij}$ on the quadratic terms.
Averaging a linear term written with measure $\varrho$ against a
flat-delta correlator of amplitude $G$ gives a quadratic term
$\tfrac12\int\varrho\,\d\yco\,(\varrho G)\hat\field^2$, so the quadratic
coefficient must be $\varrho G$, not $G$. The earlier action therefore sat
in no single convention, and its response equations carried the bracket
$\mathcal{Z}[\measure]$ while its forward equations were flat. That bracket
is a convention artefact; in the $\volel$ convention it vanishes.
\end{panel}

\subsection{Response fields and the anomalous factor of $\im$}
\label{sec:response-rotation}

Before writing the action, it is worth being explicit about a factor of
$\im$ that is otherwise easy to lose, since it fixes the sign convention
used throughout the rest of this section (Section~\ref{sec:instanton-eqs}
in particular). The standard Janssen--De~Dominicis construction introduces
response fields via the Fourier representation of the functional
delta-function enforcing the Langevin equations of motion
\eqref{eq:langevin-y},
\begin{equation}
    \delta\big[\dot\field-\mom-\advcoef\partial_\yco\field-\noise{\field}\big]
    =
    \int \Dm{\hat\field} \,
    \exp\!\left[\im\!\int_{\Ninit}^{\Nfinal}\!\!\d\Ny\int_{-1}^1\!\volel\,\d\yco\;
    \hat\field\big(\dot\field-\mom-\advcoef\partial_\yco\field-\noise{\field}\big)\right] ,
\end{equation}
and similarly for the $\mom$-equation with a genuine auxiliary field
$\hat\mom$; $\hat\field,\hat\mom$ here are real-contour fields, dual to
$\field,\mom$. Averaging over the Gaussian noise
\eqref{eq:noise-angle-averaged}, whose covariance is
$\Sigvol{ij}\,\delta(\yco-\yco')/\volel$ with
\begin{equation}
    \Sigvol{ij}(\yco,\Ny) \equiv 2\,\Dcm{ij}\atloc\,\vH(\yco,\Ny) = \volel(\yco,\Ny)\,\Dnoise{ij}(\yco,\Ny) ,
    \qquad
    \vH = \frac{4\pi}{3\,\aHloc^3} ,
    \label{eq:Sigvol-definition}
\end{equation}
produces the weight $\e{-\mathcal S}$ with
\begin{equation}
    \mathcal S
    =
    -\im\!\int\!\volel\,\d\yco\,\d\Ny\,\big(\hat\field(\dot\field-\mom-\advcoef\partial_\yco\field)
    + \hat\mom(\ldots)\big)
    +
    \tfrac12\!\int\!\volel\,\d\yco\,\d\Ny\,\big(\Sigvol{\field}\hat\field^2 + 2\Sigvol{\field\mom}\hat\field\hat\mom + \Sigvol{\mom}\hat\mom^2\big) ,
\end{equation}
which is manifestly complex for real $\hat\field,\hat\mom$: this is the
anomalous factor of $\im$. Extremizing $\mathcal S$ with respect to
$\hat\field$ (at fixed real $\field,\mom$) shows the saddle point is pure
imaginary, $\hat\field^\star\propto\im\times(\text{real})$, so the
instanton does not lie on the real $\hat\field$-contour. We therefore
rotate onto the imaginary axis,
\begin{equation}
    \hat\field \equiv \im\,\rfield , \qquad \hat\mom \equiv \im\,\rmom ,
    \label{eq:response-rotation}
\end{equation}
with $\rfield,\rmom$ real. This is the substitution that makes the
saddle-point equations real, and $\rfield,\rmom$ --- not
$\hat\field,\hat\mom$ --- are the fields appearing everywhere below.
Substituting \eqref{eq:response-rotation} into $\mathcal S$ turns the
linear term into $+\rfield(\dot\field-\mom-\advcoef\partial_\yco\field)$
(the two factors of $\im$ cancelling against the overall $-\im$) and the
quadratic term into
$-\tfrac12(\Sigvol{\field}\rfield^2+2\Sigvol{\field\mom}\rfield\rmom+\Sigvol{\mom}\rmom^2)$
(each factor of $\hat X$ contributing $\im^2=-1$), reproducing exactly the
action written below.

\subsection{The action}

With the measure $\volel$, response fields $\rfield$, $\rmom$ conjugate to
$\field$, $\mom$, the gradient prefactor
\begin{equation}
    \gpref(\yco,\Ny) \equiv \frac{1}{\rout^2\aHloc^2} = \e{-2\Deltas_{\rm loc}(\yco,\Ny)}
    \label{eq:gpref-definition}
\end{equation}
of \eqref{eq:gradient-term-y}, and the advection terms of
\eqref{eq:advection-operator} in both bracketed equations of motion, the
action is
\begin{multline}
    S
    =
    \int_{\Ninit}^{\Nfinal}\! \d\Ny \int_{-1}^1\! \volel(\yco,\Ny)\,\d\yco \; \Bigg[
    \rfield \Big( \dot\field - \mom - \advcoef\,\partial_\yco\field \Big)
    \\
    +
    \rmom \left( \dot\mom + \big(3-\epsSR\atloc\big)\mom + \frac{\Vpot'(\field)}{\Hub^2\atloc}
    - \gpref\,\Lop\field - \advcoef\,\partial_\yco\mom \right)
    -
    \frac{\Sigvol{\field}}{2}\rfield^2
    -
    \Sigvol{\field\mom}\,\rfield\,\rmom
    -
    \frac{\Sigvol{\mom}}{2}\rmom^2
    \Bigg] ,
    \label{eq:msr-action}
\end{multline}
where an overdot denotes $\partial/\partial\Ny$ at fixed $\yco$, and the
cross term $-\Sigvol{\field\mom}\rfield\rmom$ has coefficient $1$, not
$\tfrac12$, being the off-diagonal part of $-\tfrac12\rfield_i\Sigvol{ij}\rfield_j$.
The amplitudes $\Sigvol{ij}=2\Dcm{ij}\vH$ are the undiluted single-patch
amplitudes times the comoving Hubble volume; the geometric dilution lives
in the measure. Since $\Sigvol{ij}=\volel\Dnoise{ij}$, the forward
sourcing $\Sigvol{ij}\rfield_j$ equals $\Dnoise{ij}\,(\volel\rfield_j)$: the
product of the diluted coefficient with the flat-normalized response field
of the earlier draft. The response fields of this document are related to
those of the 6 July draft by $\rfield_{\rm here}=\rfield_{\rm old}/\Cpref(\Ny)$.

\begin{panel}
\semibold{Check of self-adjointness, with the drift term.} Writing $\Lop$
(dropping the $\Ny$-dependent overall constant $\e{\Deltas}\cdot4/\Deltas^2$
from \eqref{eq:Lop-definition}, which factors out of the check) as
$p(\yco)\,Q''+q(\yco)\,Q'$ with $p(\yco)=e^{\Deltas\yco}$,
$q(\yco)=-\tfrac12\Deltas e^{\Deltas\yco}$, and writing
$\weightfn(\yco)\equiv \measure(\yco,\Ny)\,p(\yco) = e^{-\Deltas\yco/2}$
(the product of the self-adjoint measure with the coefficient of $Q''$,
\emph{not} its inverse: $\measure=\weightfn^3$, and $\weightfn'=-\tfrac12\Deltas\,\weightfn$),
\begin{align*}
    \int \measure\,\d\yco\; P\,\Lop Q
    & \propto
    \int \weightfn\,\d\yco\;P\,Q'' \;-\; \frac{\Deltas}{2}\int\weightfn\,\d\yco\;P\,Q'
    \\
    & =
    \Big[\weightfn P Q'\Big]_{\mathrm{bdry}} - \int \weightfn'\,P\,Q'\,\d\yco - \int\weightfn\,P'\,Q'\,\d\yco - \frac{\Deltas}{2}\int\weightfn\,P\,Q'\,\d\yco
    \\
    & =
    \Big[\weightfn P Q'\Big]_{\mathrm{bdry}} - \int \weightfn\, P'\,Q'\,\d\yco ,
\end{align*}
using $\weightfn'=-\tfrac12\Deltas\weightfn$ to cancel the two
$\weightfn'PQ'$-type terms exactly. The factor $\Cpref(\Ny)$ relating
$\measure$ to $\volel$ is $\yco$-independent and multiplies both sides, so
the check is identical in the $\volel$ convention. The bulk term
$-\int\weightfn P'Q'\,\d\yco$ is manifestly symmetric in $P,Q$;
self-adjointness therefore requires only that the boundary term
$\big[\weightfn(PQ'-QP')\big]_{\mathrm{bdry}}$ vanish.

In the action the operator acts on $\field$ with the multiplier
$\gpref\rmom$, so the relevant pair is $Q=\delta\field$, $P=\gpref\rmom$.
At $\yco=+1$ the forward condition $\partial_\yco\field=0$ gives $Q'=0$,
and the response-sector condition must make $P'=\partial_\yco(\gpref\rmom)$
vanish there (Section~\ref{sec:bcs}). These are one admissible condition
per sector, matching the characteristic count of Section~\ref{sec:bcs}
(one incoming characteristic at the core in each sector). The response
condition supplies no data: the boundary term of the action is affine in
$(\field,\mom)$ with coefficients linear in the response fields, so
varying it with respect to $(\field,\mom)$ gives homogeneous conditions
and any target values differentiate away. At $\yco=-1$,
$\rfield=\rmom=0$ (no response where the outer data is pinned) gives $P=0$.
\end{panel}
% ======================================================================
\section{The Fokker--Planck Hamiltonian}
\label{sec:hfp}
% ======================================================================

The MSR action \eqref{eq:msr-action} has the form $\int\d\Ny\,[p\dot q-H]$.
The Hamiltonian it defines generates both the forward and the response
equations, and it is the object from which the instanton equations should
be obtained. This section defines it, verifies that it reproduces the
equations of Section~\ref{sec:instanton-eqs}, records its conservation
properties, and gives its closed form for the gradient-free instanton.

\subsection{Definition and canonical variables}

Write $\langle\cdot\rangle \equiv \int_{-1}^{1}\volel(\yco,\Ny)\,\d\yco\,(\cdot)$.
Grouping \eqref{eq:msr-action} by whether a term multiplies an
$\Ny$-derivative,
\begin{equation}
    S = \int_{\Ninit}^{\Nfinal}\d\Ny\,\Big[\big\langle\rfield\dot\field+\rmom\dot\mom\big\rangle - \Hfp\Big] ,
    \qquad
    \Hfp = \Big\langle \rfield\,b_\field + \rmom\,b_\mom
    + \tfrac12\Sigvol{\field}\rfield^2 + \Sigvol{\field\mom}\rfield\rmom + \tfrac12\Sigvol{\mom}\rmom^2 \Big\rangle ,
    \label{eq:hfp-definition}
\end{equation}
where $b_\field$, $b_\mom$ are the drift terms of the Langevin system
\eqref{eq:langevin-y},
\begin{equation}
    b_\field = \mom + \advcoef\,\partial_\yco\field ,
    \qquad
    b_\mom = -\big(3-\epsSR\atloc\big)\mom - \frac{\Vpot'(\field)}{\Hub^2\atloc}
    + \gpref\,\Lop\field + \advcoef\,\partial_\yco\mom .
    \label{eq:drift-definition}
\end{equation}
This is the Freidlin--Wentzell Hamiltonian
$H(q,\theta)=\langle\theta,b(q)\rangle+\tfrac12\langle\theta,a\theta\rangle$
of large-deviation theory, with $\theta=(\rfield,\rmom)$ and
$a=\Sigvol{ij}$. The sign bookkeeping is that the $\rmom$ bracket in
\eqref{eq:msr-action} is $-b_\mom$ and enters $S$ with a plus sign, while
the quadratic terms enter with a minus sign; both flip on passing to
$-\Hfp$.

\begin{panel}
\semibold{The measure is part of the symplectic form.} The symplectic
form is $\int\d\yco\,\volel\,\rfield\,\delta\field$, so the momentum
density conjugate to $\field$ with respect to the flat measure $\d\yco$ is
$\pfield=\volel\rfield$, and likewise $\pmom=\volel\rmom$, not the response
fields themselves. In these variables the drift terms are unchanged (the
measure cancels against the density) and the quadratic terms pick up
$1/\volel$:
\begin{equation}
    \Hfp = \int_{-1}^{1}\d\yco\,\Big[\pfield\,b_\field + \pmom\,b_\mom
    + \tfrac12\Dnoise{\field}\pfield^2 + \Dnoise{\field\mom}\pfield\pmom + \tfrac12\Dnoise{\mom}\pmom^2\Big] ,
    \label{eq:hfp-canonical}
\end{equation}
using $\Sigvol{ij}/\volel=\Dnoise{ij}$. The canonical form carries the
diluted coefficients $\Dnoise{ij}=2\Dcm{ij}/\ncount$ of
\eqref{eq:Dnoise-diag} with no measure factor at all. This is the
convention independence of Section~\ref{sec:measure-choice} seen from the
Hamiltonian side: every convention gives the same $\Hfp$ in canonical
variables. The time derivative of $\volel$ that appears when the response
equations are written for $\rfield$ rather than $\pfield$ is the only
place the convention shows.
\end{panel}

\subsection{Hamilton's equations}

\begin{calculation}
\semibold{Forward sector.} From \eqref{eq:hfp-canonical},
\begin{equation}
    \dot\field = \frac{\delta\Hfp}{\delta\pfield}
    = b_\field + \Dnoise{\field}\pfield + \Dnoise{\field\mom}\pmom
    = \mom + \advcoef\,\partial_\yco\field + \Sigvol{\field}\rfield + \Sigvol{\field\mom}\rmom ,
\end{equation}
which is \eqref{eq:inst-phi}; likewise $\delta\Hfp/\delta\pmom$ gives
\eqref{eq:inst-pi}. These are exact: no coefficient is frozen.

\semibold{Response sector.} $\dot\pfield=-\delta\Hfp/\delta\field$. Varying
at frozen coefficients (the full variation is discussed in
Section~\ref{sec:instanton-eqs}), term by term:
\begin{itemize}
    \item $\int\pfield\,\advcoef\,\partial_\yco\field$ contributes
    $-\partial_\yco(\advcoef\,\pfield)$, after one integration by parts
    in $\yco$;
    \item $\int\pmom\,(-\Vpot'/\Hub^2)$ contributes $-\pmom\,\Vpot''/\Hub^2$;
    \item $\int\pmom\,\gpref\,\Lop\field = \int\volel\,(\gpref\rmom)\,\Lop\field$
    contributes, by self-adjointness of $\Lop$ with respect to
    $\volel\,\d\yco$, $\volel\,\Lop(\gpref\rmom)$, with the prefactor
    inside the operator.
\end{itemize}
Hence
\begin{equation}
    \dot\pfield = \partial_\yco(\advcoef\,\pfield) + \pmom\frac{\Vpot''}{\Hub^2} - \volel\,\Lop(\gpref\rmom) .
\end{equation}
Substituting $\pfield=\volel\rfield$, expanding
$\dot\pfield=\dot\volel\rfield+\volel\dot{\rfield}$, and dividing by
$\volel$,
\begin{equation}
    \dot{\rfield}
    = \Big[-\partial_\Ny\ln\volel + \partial_\yco\advcoef + \advcoef\,\partial_\yco\ln\volel\Big]\rfield
    + \advcoef\,\partial_\yco\rfield + \frac{\Vpot''}{\Hub^2}\rmom - \Lop(\gpref\rmom)
    = \advcoef\,\partial_\yco\rfield + \frac{\Vpot''}{\Hub^2}\rmom - \Lop(\gpref\rmom) ,
\end{equation}
the bracket being $\mathcal{Z}[\volel]=0$ of \eqref{eq:Z-conventions};
the cancellation is displayed in Section~\ref{sec:instanton-eqs}. The
$\rmom$ equation follows in the same way from $\dot\pmom=-\delta\Hfp/\delta\mom$.
\end{calculation}

\subsection{Reduction to one dimension and the Legendre check}
\label{sec:hfp-1d}

In the gradient-free limit ($\advcoef=0$, $\Lop=0$, a single patch with
$\ncount=1$ and $\volel=1$), with the conventions of the single-trajectory
instanton ($\Dnoise{\field}=2\Dcm{11}$, response fields $P_1,P_2$),
\begin{equation}
    \Hfp = P_1\,\mom + P_2\Big[-(3-\epsSR)\mom - \frac{\Vpot'}{\Hub^2}\Big]
    + \Dcm{11}P_1^2 + 2\Dcm{12}P_1P_2 + \Dcm{22}P_2^2 .
    \label{eq:hfp-1d}
\end{equation}
The Legendre check fixes every sign at once: with
$\dot\field=\mom+2\Dcm{11}P_1+2\Dcm{12}P_2$ and
$\dot\mom=b_\mom+2\Dcm{12}P_1+2\Dcm{22}P_2$,
\begin{equation}
    P_1\dot\field + P_2\dot\mom - \Hfp = \Dcm{11}P_1^2 + 2\Dcm{12}P_1P_2 + \Dcm{22}P_2^2 ,
\end{equation}
which is the on-shell action density $S=\int\Dcm{11}P_1^2\,\d\Ny$ of the
single-trajectory instanton when $\Dcm{12}=\Dcm{22}=0$. In the onion model
the same identity reads
$\langle\rfield\dot\field+\rmom\dot\mom\rangle-\Hfp
=\langle\tfrac12\Sigvol{\field}\rfield^2+\Sigvol{\field\mom}\rfield\rmom+\tfrac12\Sigvol{\mom}\rmom^2\rangle$,
the on-shell action density. This settles a factor-of-two question raised
on 6--7 July 2026 between $\tfrac12\Dnoise{\field}\rfield^2$ and
$\Dcm{11}P_1^2$: with $\Dnoise{\field}=2\Dcm{11}/\ncount$ the two agree
at $\ncount=1$.

\subsection{Conservation, non-conservation, and the balance law}

In the gradient-free limit $\Hfp$ has no explicit $\Ny$-dependence and no
spatial boundary, so it is exactly conserved along the instanton. In the
onion model it is not, for three independent reasons.

\para{Explicit $\Ny$-dependence} $\Deltas$, $\Deltas_{\rm loc}$ and
$\ncount$ depend on $a(\Ny)$ through \eqref{eq:Deltas-closed}; at fixed
fields $\partial\Deltas/\partial\Ny=1$. The expansion is an explicit clock
and $\partial_\Ny\Hfp|_{\rm expl}\neq0$.

\para{Boundary flux} The bulk cancellation in $\d\Hfp/\d\Ny$ leaves the
bilinear boundary form that the variational derivation discards,
\begin{equation}
    \Phi = \Big[\volel\advcoef\big(\rfield\dot\field+\rmom\dot\mom\big)
    + \Cpref\kappa\,\weightfn\big(\gpref\rmom\,\partial_\yco\dot\field - \dot\field\,\partial_\yco(\gpref\rmom)\big)\Big]_{\yco=-1}^{\yco=+1} ,
    \label{eq:hfp-flux}
\end{equation}
with $\kappa=\e{\Deltas}\cdot4/\Deltas^2$ the constant factored out of
$\Lop$ and $\weightfn=\e{-\Deltas\yco/2}$ as in the self-adjointness
panel. At $\yco=+1$ the gradient piece vanishes under the boundary
conditions of Section~\ref{sec:bcs} ($\partial_\yco\field=0$ for all $\Ny$
implies $\partial_\yco\dot\field=0$, and the response condition kills
$\partial_\yco(\gpref\rmom)$), but the advective piece survives because
$\advcoef(1)\neq0$: the core is a moving boundary, a piston sweeping into
the shrinking horizon. At $\yco=-1$, $\advcoef(-1)=0$ and $\rmom(-1)=0$,
but the Dirichlet datum $\phinl(\Ny)$ is time dependent and
$\partial_\yco(\gpref\rmom)$ is unconstrained there, so the second gradient
term survives: the outer edge is a driven boundary. Neither flux can be
removed by a choice of boundary condition; both are physical.

\para{Truncation} If the response equations actually solved are not the
exact variation of $\Hfp$, write them as $-\delta\Hfp/\delta\field+R_\field$
and $-\delta\Hfp/\delta\mom+R_\mom$. The residual $R$ then contributes a
bulk drift. The exact balance law is
\begin{equation}
    \frac{\d\Hfp}{\d\Ny}
    = \frac{\partial\Hfp}{\partial\Ny}\bigg|_{\rm expl} + \Phi
    + \big\langle\dot\field\,R_\field + \dot\mom\,R_\mom\big\rangle .
    \label{eq:balance-law}
\end{equation}
This, not $\d\Hfp/\d\Ny=0$, is the conservation statement for the onion
model. It separates three failure modes that would otherwise be conflated
and is a pointwise-in-$\Ny$ consistency test. When the response sector is
generated from $\Hfp$ itself (Section~\ref{sec:collocation-scheme}),
$R=0$ by construction. With a truncated response sector the solution does
not extremize $S$, and the on-shell action then carries an error linear in
$R$ rather than quadratic.

\subsection{$\Hfp$ is not zero; the Maupertuis relation}

\begin{panel}
\semibold{The vanishing-Hamiltonian expectation, and why it fails.} In the
noiseless phase every term of $\Hfp$ carries a response field, so
$\Hfp=0$. One might expect a smooth continuation from the noiseless phase
to enforce $\Hfp=0$ throughout, as an unenforced constraint. It does not.
$H=0$ holds for a trajectory emanating from a fixed point of the noiseless
flow with vanishing response fields, which forces an infinite transition
time. Neither property holds here: $(\phiinit,\mom_{\rm SR,init})$ is not a
fixed point, and the duration $\DN+\dNstar$ is finite and prescribed. The
response fields therefore do not grow smoothly from zero at $\Ninit$; they
jump, from zero before $\Ninit$ to a value fixed by backward integration
from the terminal source. The jump is the shadow of the modelling choice
that noise acts only on $[\Ninit,\Nfinal]$, and $\Hfp$ is the shadow price
of that constraint. It is the same structural fact as the largest-time
question of Section~\ref{sec:open-issues}, seen from the other end of the
interval. $\Hfp=0$ is the condition that would determine $\dNstar$ if the
duration were free; imposing it on top of the present boundary value
problem would over-determine it.
\end{panel}

For the gradient-free instanton $\Hfp$ is constant, and the envelope
theorem for a fixed-endpoint, fixed-duration problem gives
$\partial S/\partial T=-\Hfp$. Varying $\dNstar$ at fixed $(\Nfinal,\DN)$
holds both endpoints fixed and varies only the duration, so
\begin{equation}
    \Hfp = -\frac{\partial S}{\partial\dNstar}\bigg|_{\Nfinal,\DN} .
    \label{eq:maupertuis}
\end{equation}
Since $S$ increases with $\dNstar$, $\Hfp<0$. In the onion model the
relation fixes $\Hfp$ at the terminal time rather than a constant, and the
profile $\Hfp(\Ny)-\Hfp(\Nfinal)$ is the accumulated non-conservation of
\eqref{eq:balance-law}.

\para{Closed form, gradient-free case} The terminal conditions
$P_2(\Ntotal)=0$, $P_1(\Ntotal)=\lag$ pin the conserved value of
\eqref{eq:hfp-1d}:
\begin{equation}
    \Hfp = \lag\,\mom(\Ntotal) + \Dcm{11}(\Ntotal)\,\lag^2 ,
    \qquad\text{hence}\qquad
    \frac{\partial S}{\partial\dNstar}\bigg|_{\Nfinal,\DN} = -\Hfp ,
    \label{eq:hfp-closed-form}
\end{equation}
with $\Ntotal=\DN+\dNstar$. Every quantity on the right is available from
a converged solution, and the left is a finite difference across a
$\dNstar$ grid. The identity holds despite the free $\mom$ endpoint
because $P_2(\Ntotal)=0$ removes the corresponding boundary variation. It
is an over-determined check: $\Hfp$ should be constant in $\Ny$ and equal
to the finite difference, and the same residual $R$ controls the failure
of both.

\begin{panel}
\semibold{Response fields as conjugate momenta, not as an adjoint
transport.} With the drift $b(u)$ and noise amplitude $\Sigvol{}(u)$
depending on the state $u=(\field,\mom)$, Hamilton's equations for
$\mathcal{H}(u,P)=P\cdot b(u)+\tfrac12P\cdot\Sigvol{}(u)\cdot P$ give
\begin{equation}
    \dot P_i = -P_j\,\frac{\partial b_j}{\partial u_i} - \frac12 P_j\,\frac{\partial\Sigvol{jk}}{\partial u_i}\,P_k .
\end{equation}
Only when $\partial\Sigvol{}/\partial u=0$ does $P$ transport as the
adjoint tangent map of the drift, $\dot P=-(\partial b/\partial u)^{\transpose}P$.
Otherwise no adjoint-transport statement survives: transposing the full
forward Jacobian $\partial f/\partial u$ with $f=b+\Sigvol{}P$ gives the
quadratic term with coefficient $1$ instead of $\tfrac12$. The object to
implement is the scalar $\mathcal{H}$; both sectors are its Hamilton
equations, and the factor of $\tfrac12$ is automatic. The full Jacobian
of the joint flow is $\Omega\,\mathrm{Hess}\,\mathcal{H}$ with
$\Omega=\bigl(\begin{smallmatrix}0&I\\-I&0\end{smallmatrix}\bigr)$, so the
structural check on any implementation is that $\Omega J_{\rm full}$ is
symmetric: three block identities, replacing the earlier test
$J_{\rm resp}+J_{\rm fwd}^{\transpose}=0$, which is valid only when
$\partial\Sigvol{}/\partial u=0$ and then only against the full forward
Jacobian including $\partial_u(\Sigvol{}P)$. For the paper, the response
fields should be described as conjugate momenta of $\Hfp$; the
adjoint-transport language is safe only under a stated
$\partial\Sigvol{}/\partial u=0$ assumption. The decision taken on
8 October 2026 is that both sectors are obtained by differentiating a
single discrete $\Hfp$ (Section~\ref{sec:collocation-scheme}); the
continuum response equations displayed in Section~\ref{sec:instanton-eqs}
are for understanding, and are not transcribed.
\end{panel}
% ======================================================================
\section{Instanton equations}
\label{sec:instanton-eqs}
% ======================================================================

The saddle point of \eqref{eq:msr-action} is given by Hamilton's equations
of $\Hfp$ (Section~\ref{sec:hfp}). The forward equations are exact and
short. The response equations are the full variation of $\Hfp$ with
respect to $\field$ and $\mom$, and in the onion model that variation
includes the state dependence of the coordinate scaffold through the core
node. This section first gives the response equations of the gradient-free
instanton in full, as intuition-building material, then the leading terms
of the onion response equations, and finally the statement of what the
full variation contains.

\subsection{The gradient-free instanton}
\label{sec:fullinstanton-eqs}

With the Hamiltonian \eqref{eq:hfp-1d}, writing $u=(\field,\mom)$,
$f=(\mom,\,-(3-\epsSR)\mom-\Vpot'/\Hub^2)$ and treating
$\Hub^2(\field,\mom)$, $\epsSR(\field,\mom)$, $\Dcm{ij}(\field,\mom)$ as the
local functions they are, the forward equations are $\dot u=\partial\Hfp/\partial P$,
\begin{equation}
    \dot\field = \mom + 2\Dcm{11}P_1 + 2\Dcm{12}P_2 ,
    \qquad
    \dot\mom = -(3-\epsSR)\mom - \frac{\Vpot'}{\Hub^2} + 2\Dcm{12}P_1 + 2\Dcm{22}P_2 ,
    \label{eq:fi-forward}
\end{equation}
and the response equations are $\dot P=-\partial\Hfp/\partial u$,
\begin{subequations}
\label{eq:fi-response}
\begin{align}
    \dot P_1
    & = P_2\left[\frac{\Vpot''}{\Hub^2} - \frac{\partial\epsSR}{\partial\field}\,\mom
    - \frac{\Vpot'}{\Hub^4}\frac{\partial\Hub^2}{\partial\field}\right]
    - \left[\frac{\partial\Dcm{11}}{\partial\field}P_1^2 + 2\frac{\partial\Dcm{12}}{\partial\field}P_1P_2 + \frac{\partial\Dcm{22}}{\partial\field}P_2^2\right] ,
    \label{eq:fi-P1}
    \\
    \dot P_2
    & = -P_1 + (3-\epsSR)P_2
    - P_2\left[\frac{\partial\epsSR}{\partial\mom}\,\mom + \frac{\Vpot'}{\Hub^4}\frac{\partial\Hub^2}{\partial\mom}\right]
    - \left[\frac{\partial\Dcm{11}}{\partial\mom}P_1^2 + 2\frac{\partial\Dcm{12}}{\partial\mom}P_1P_2 + \frac{\partial\Dcm{22}}{\partial\mom}P_2^2\right] .
    \label{eq:fi-P2}
\end{align}
\end{subequations}
For canonical kinetics in reduced Planck units, $\epsSR=\mom^2/2$ and
$\Hub^2=\Vpot/(3-\epsSR)$, so
\begin{equation}
    \frac{\partial\epsSR}{\partial\field}=0 , \quad
    \frac{\partial\epsSR}{\partial\mom}=\mom , \quad
    \frac{\partial\Hub^2}{\partial\field}=\Hub^2\frac{\Vpot'}{\Vpot} , \quad
    \frac{\partial\Hub^2}{\partial\mom}=\frac{\Hub^2\mom}{3-\epsSR} ,
\end{equation}
and the bracketed drift terms become
\begin{equation}
    \dot P_1 \supset P_2\,\frac{\Vpot''-\Vpot'^2/\Vpot}{\Hub^2} = P_2\,(3-\epsSR)\,(\ln\Vpot)'' ,
    \qquad
    \dot P_2 \supset -P_2\Big[2\epsSR + \frac{\Vpot'}{\Vpot}\mom\Big] .
    \label{eq:fi-canonical-terms}
\end{equation}
For the massless de Sitter diffusion model $\Dcm{11}=\Hub^2/8\pi^2$,
$\Dcm{12}=\Dcm{22}=0$, the diffusion terms are
$-(\partial_\field\Dcm{11})P_1^2 = -\Dcm{11}(\Vpot'/\Vpot)P_1^2$ and
$-(\partial_\mom\Dcm{11})P_1^2 = -\Dcm{11}\,\mom P_1^2/(3-\epsSR)$. These
are quadratic in $P_1$ and make the backward pass a Riccati flow; see the
end of this section.

\begin{calculation}
\semibold{The dropped terms are not small.} The term retained in the
earlier response sector was $P_2\Vpot''/\Hub^2$. The exact coefficient of
$P_2$ in $\dot P_1$ is
$-\partial b_\mom/\partial\field = (3-\epsSR)(\ln\Vpot)''$, because
$\Vpot'/\Hub^2=(3-\epsSR)\Vpot'/\Vpot$. For $\Vpot\propto\field^p$,
\begin{equation}
    (\ln\Vpot)'' = -\frac{p}{\field^2} ,
    \qquad
    \frac{\Vpot''}{\Vpot} = \frac{p(p-1)}{\field^2} ,
    \qquad
    \frac{\Vpot(\ln\Vpot)''}{\Vpot''} = -\frac{1}{p-1} ,
\end{equation}
so the exact drift term is $-1/(p-1)$ times the retained one. For a
quadratic potential the two are equal and opposite: the retained term has
the wrong sign at the same magnitude. For a quartic potential the factor
is $-1/3$. For an ultra-slow-roll plateau both vanish, which is a null
test. In slow-roll language the omitted piece is
$-P_2\Vpot'^2/(\Vpot\Hub^2)\approx-6\epsSR P_2$ against the retained
$\approx3\eta_V P_2$: the same order. The $\dot P_2$ corrections are
genuinely $\Or(\epsSR)$.
\end{calculation}

\begin{panel}
\semibold{Terms previously dropped.} The 6 July draft stated that the
derivatives $\partial\Hub^2/\partial\field$, $\partial\epsSR/\partial\field$,
$\partial\Dcm{ij}/\partial\field$ and the analogous $\mom$ derivatives
could be dropped as ``a higher-order effect in the same sense as other
terms already truncated in the gradient expansion''. That was wrong: the
calculation above shows the omission is leading order for the drift, and
the omitted diffusion terms change the character of the backward pass.
Every action value computed with the truncated response sector before
October 2026 is contaminated at that order and must be recomputed; the
scaling fits in $\dNstar$ and $\DN$, the threshold locus and the comparison
with the Vennin exponent are all affected.

Physics of the derivative terms. $\Hub$ and $\epsSR$ are not independent
variables of the original Langevin equations; they come from the
Einstein constraint evaluated on the local state. The MSR action is the
action of the reduced stochastic system on $(\field,\mom)$ obtained after
that constraint has been used. The derivative terms therefore appear only
in the response sector; the Langevin equations and the forward sector are
unchanged. A noise realization moves $\Hub$, and the response fields must
know it. For the onion model, the local $\Hub^2$ used in each shell is a
truncated constraint (no gradient or curvature terms); differentiating it
is consistent with the separate-universe closure of the model and is the
convention adopted here.

A further dropped term lies in the model rather than in the response
sector. Constant-$\Ny$ hypersurfaces and constant-proper-time
hypersurfaces are tilted relative to one another across shells, a
shift-vector-type effect identified on 2 July 2026 and dropped at the same
order as the $\partial\Dcm{ij}/\partial\field$ feedback. The advection
coefficient $\advcoef$ is a shift vector in the optical metric of
Section~\ref{sec:bcs}; the tilt is the piece of that shift not captured by
the comoving relabelling.
\end{panel}

\subsection{The onion model}
\label{sec:onion-eqs}

The forward equations are, exactly,
\begin{subequations}
\label{eq:instanton-eqs}
\begin{align}
    \dot\field
    & =
    \mom + \advcoef\,\partial_\yco\field + \Sigvol{\field}\,\rfield + \Sigvol{\field\mom}\,\rmom ,
    \label{eq:inst-phi}
    \\
    \dot\mom
    & =
    -\big(3-\epsSR\atloc\big)\mom - \frac{\Vpot'(\field)}{\Hub^2\atloc}
    + \gpref\,\Lop\field
    + \advcoef\,\partial_\yco\mom
    + \Sigvol{\mom}\,\rmom + \Sigvol{\field\mom}\,\rfield ,
    \label{eq:inst-pi}
\end{align}
with $\advcoef$ built from the chain-rule rate \eqref{eq:Deltas-dot-closed}.
The response equations, at frozen coefficients, are
\begin{align}
    \dot{\rfield}
    & =
    \advcoef\,\partial_\yco\rfield
    + \frac{\Vpot''(\field)}{\Hub^2\atloc}\rmom
    - \Lop\big(\gpref\,\rmom\big) ,
    \label{eq:inst-rphi}
    \\
    \dot{\rmom}
    & =
    \advcoef\,\partial_\yco\rmom
    - \rfield + \big(3-\epsSR\atloc\big)\rmom .
    \label{eq:inst-rpi}
\end{align}
\end{subequations}
The response fields take the place of the noises in
\eqref{eq:inst-phi}--\eqref{eq:inst-pi}, with the cross term
$\Sigvol{\field\mom}$ coupling each response field to the other noise
channel. There is no multiplicative term in the response equations: the
coefficient $\mathcal{Z}[\volel]$ of \eqref{eq:Z-conventions} vanishes.
In the gradient-free limit these reduce to the leading terms of
\eqref{eq:fi-response}.

\begin{calculation}
\semibold{The three-way cancellation $\mathcal{Z}[\volel]=0$.} Varying
\eqref{eq:msr-action} with respect to $\field$ and dividing by $\volel$,
the terms multiplying $\rfield$ without a $\yco$-derivative come from the
$\Ny$ integration by parts of $\volel\rfield\dot\field$ and the $\yco$
integration by parts of $-\volel\rfield\advcoef\partial_\yco\field$:
\begin{equation}
    \mathcal{Z}[\volel] = \underbrace{-\partial_\Ny\ln\volel}_{\text{measure, time}}
    + \underbrace{\advcoef\,\partial_\yco\ln\volel}_{\text{measure, space}}
    + \underbrace{\partial_\yco\advcoef}_{\text{coefficient}} .
\end{equation}
With $\volel=2\pi\rout^3\Deltas\,\e{-\frac32(\yco+1)\Deltas}$ and
$\advcoef=(\yco+1)\dot{\Deltas}/\Deltas$,
\begin{equation}
    -\partial_\Ny\ln\volel = -\frac{\dot{\Deltas}}{\Deltas} + \frac32(\yco+1)\dot{\Deltas} ,
    \qquad
    \advcoef\,\partial_\yco\ln\volel = -\frac32(\yco+1)\dot{\Deltas} ,
    \qquad
    \partial_\yco\advcoef = \frac{\dot{\Deltas}}{\Deltas} ,
\end{equation}
and the sum vanishes identically. The mechanism is three-way, not two-way.
The $\yco$-dependent parts, $\pm\tfrac32(\yco+1)\dot{\Deltas}$, are the
time and space images of the exponential factor and cancel between the
two measure terms. The prefactor $\Deltas(\Ny)$ is invisible to
$\partial_\yco\ln\volel$; its time derivative $-\dot{\Deltas}/\Deltas$ is
cancelled by the derivative of the advection coefficient, not by a
measure term. With $\measure$ in place of $\volel$ the prefactor is absent
and $\partial_\yco\advcoef$ has nothing to cancel against: that leftover is
the $1/\Deltas$ in the old bracket. The cancellation is a continuity
equation, $\partial_\Ny\volel=\partial_\yco(\volel\advcoef)$: $\advcoef$ is
minus the coordinate velocity, and only the true Jacobian of the map is
transported by the coordinate flow. The same cancellation holds for any
state-dependent $\Deltas(\Ny)$ as long as $\advcoef$ and $\volel$ are built
from the same one. The $\Lop$ term is untouched, and the $\rmom$ row goes
through identically.
\end{calculation}

\para{Operator ordering in the gradient term} The response gradient term
is $\Lop(\gpref\rmom)$, with the prefactor $\gpref=\e{-2\Deltas_{\rm loc}}$
inside the operator. The action term is $-\int\volel\,\rmom\,\gpref\,\Lop\field$,
and self-adjointness of $\Lop$ moves the operator onto the whole product
$\gpref\rmom$. Writing $\Lop=\kappa[p\,\partial_\yco^2+q\,\partial_\yco]$,
\begin{equation}
    \Lop(\gpref\rmom) - \gpref\,\Lop\rmom
    = 2\kappa\,p\,\gpref'\,\partial_\yco\rmom + \kappa\big(p\,\gpref''+q\,\gpref'\big)\rmom ,
    \label{eq:ordering-terms}
\end{equation}
both of which are missing if the prefactor is placed outside. In
covariant form, $\nabla^2(\gpref P)=\gpref\nabla^2P+2\nabla\gpref\cdot\nabla P+P\nabla^2\gpref$:
the operator $\gpref\nabla^2$ is not self-adjoint with respect to
$\rad^2\,\d\rad$, and its adjoint is $\nabla^2(\gpref\,\cdot)$. Freezing
$\gpref$ under the variation removes $\partial\gpref/\partial\field\,\delta\field$
terms; it does not make $\gpref$ constant in $\yco$, because $\field$
varies across shells, which is the premise of the model. All three terms
are second order in gradients, so this is an inconsistency at fixed order,
not a truncation. The relative size of the missing terms is set by
$\gpref'/\gpref=-\partial_\yco\ln\Hub^2$, the shell-to-shell variation of
$\Hub$, which is $\Or(1)$ in the large-$\dNstar$ regime. By the chain rule
$\gpref'=(\partial\gpref/\partial\field)\partial_\yco\field+(\partial\gpref/\partial\mom)\partial_\yco\mom$,
so $\partial\Hub^2/\partial\field$ and $\partial\Hub^2/\partial\mom$ are
needed for a second, independent reason.

\para{The full variation} Equations \eqref{eq:inst-rphi}--\eqref{eq:inst-rpi}
are the leading terms of $-\delta\Hfp/\delta\field$,
$-\delta\Hfp/\delta\mom$. The full variation adds: the local derivative
terms of \eqref{eq:fi-response} ($\partial\Hub^2/\partial u$,
$\partial\epsSR/\partial u$, $\partial\Sigvol{ij}/\partial u$ through
$\Dcm{ij}$ and $\vH$); and, through the core node, the derivatives of
$\Deltas$, $\dot{\Deltas}$, $\advcoef$, $\volel$ and $\ncount$ with respect
to $\field(1,\Ny)$, $\mom(1,\Ny)$ and the core values of
$\partial_\yco\field$, $\partial_\yco\mom$ and the Laplacian that enter
\eqref{eq:Deltas-dot-closed}. The latter are rank-one couplings from the
whole grid into the core rows, with a quotient rule through $1/(1-\Kcoef)$.
They are generated by differentiating the discrete Hamiltonian
(Section~\ref{sec:collocation-scheme}) and are not written out. Varying
at fixed $\Deltas$ and then substituting $\Deltas[u]$ is not the same as
substituting and then varying; the anchor $\rho\equiv1$ of
Section~\ref{sec:coordinate} demands the second.

\begin{panel}
\semibold{Historical: the direct re-derivation of 6 July 2026.} The
response equations were first obtained by varying the action written with
measure $\measure$. That derivation corrected the sign of every term of
an earlier pass and found a previously missing measure-friction term
$-(\dot\measure/\measure)X$, with $-\dot\measure/\measure=\tfrac32\dot{\Deltas}\yco$.
Combined with the expansion
$\measure^{-1}\partial_\yco(\measure\advcoef X)=\advcoef\partial_\yco X+X\dot{\Deltas}[1/\Deltas-\tfrac32(\yco+1)]$,
the $\yco$-dependence cancelled and left the bracket
$X\,\dot{\Deltas}[1/\Deltas-\tfrac32]$, a single scalar per $\Ny$, on both
response fields. The sign corrections stand. The bracket is
$\mathcal{Z}[\measure]$ of \eqref{eq:Z-conventions}: an artefact of the
$\measure$ convention, removed by the physical measure, and inconsistent
with a forward sector written with the flat diffusion coefficient. The
same derivation placed the gradient prefactor outside the operator, which
is corrected above.
\end{panel}

\begin{todo}
\semibold{Riccati structure of the response sector.} With
$\partial\Sigvol{}/\partial u\neq0$ the response equations are quadratic in
the response fields, so the backward pass is a Riccati flow. The exact
linearity $\rfield=\lag\hat{\rfield}$ used to control the dynamic range of
the response fields is lost, together with the one-multiply reinstatement
of $\lag$. A per-field substitution $P=M_1/M_2$ does not linearize the
system, because the linear coupling between $P_1$ and $P_2$ defeats an
independent ratio ansatz; the correct object is the matrix Riccati
equation linearized by the $4\times4$ fundamental matrix of the
$(\field,P)$ Hamiltonian system, in which a finite-$\Ny$ caustic appears
as $M_2\to0$, a coordinate singularity rather than a step failure. The
terminal condition must be mapped into the linearized variables. The
forward sector, the Picard structure, the action form and both
conservation identities are unchanged. Whether to integrate the Riccati
form directly or use the symplectic block linearization is to be decided
on the gradient-free instanton, which is the right pilot. The linearity
itself was only ever a tree-level property: loop corrections from a
Schwinger--Keldysh effective action break it independently. See
\texttt{RECONSTRUCTION.md}, Open~4.
\end{todo}
% ======================================================================
\section{Boundary conditions}
\label{sec:bcs}
% ======================================================================

\subsection{Spatial boundary conditions}

\para{At $\yco=-1$ (fixed outer edge)} The field is pinned to the known
noiseless trajectory,
\begin{equation}
    \field(-1,\Ny) = \phinl(\Ny) ,
    \qquad
    \mom(-1,\Ny) = \pinl(\Ny) ,
    \label{eq:bc-outer}
\end{equation}
and correspondingly $\rfield(-1,\Ny) = \rmom(-1,\Ny) = 0$ (no response
where the field is pinned). In each sector this is one condition plus an
identity. Since $\advcoef(-1)=0$ and the response fields vanish there,
\eqref{eq:inst-phi} at $\yco=-1$ reads $\dot\field(-1)=\mom(-1)$, so
pinning $\field(-1)=\phinl(\Ny)$ for all $\Ny$ already implies
$\mom(-1)=\dot\phinl=\pinl$. Likewise $\rmom(-1)\equiv0$ gives
$\dot{\rmom}(-1)=0$, and \eqref{eq:inst-rpi} then reads $0=-\rfield(-1)$:
$\rfield(-1)=0$ is an identity given $\rmom(-1)=0$. The characteristic
count below allows exactly one condition at this edge, and this is it.

\para{At $\yco=+1$ (current horizon)} We impose
\begin{equation}
    \partial_\yco \field \big|_{\yco=1} = 0 ,
    \label{eq:bc-core}
\end{equation}
and in the response sector the homogeneous condition that makes the
boundary term of the action vanish, which at leading order is
$\partial_\yco\rmom|_{\yco=1}=0$ (see below). The physical content of
\eqref{eq:bc-core} is that the excised sub-horizon interior carries no
independent structure of its own, so the field attaches smoothly to
whatever lies beyond the domain, with no preferred direction singled out
at the point of attachment. $\field(1,\Ny)$ itself is left free by this
condition and is an output of the solution, not an input. The remainder
of this subsection justifies the count of one condition per sector and
the statement that no data enters at the core.

\subsection{Characteristic structure of the forward sector}
\label{sec:characteristics}

The forward sector is hyperbolic, not parabolic: the Laplacian sits in the
equation for $\dot\mom$, hence at the level of $\ddot\field$, as for the
Klein--Gordon equation it descends from. Keep the principal
(highest-derivative) parts of \eqref{eq:inst-phi}--\eqref{eq:inst-pi} with
frozen coefficients. The noise terms carry no $\yco$-derivative and do not
enter. With
\begin{equation}
    \wspeed^2(\yco,\Ny) \equiv \gpref\,\e{\Deltas}\e{\Deltas\yco}\,\frac{4}{\Deltas^2} ,
    \label{eq:wave-speed}
\end{equation}
the coefficient of $\partial_\yco^2\field$ in $\dot\mom$, the principal
system is $\dot\field=\mom+\advcoef\partial_\yco\field$,
$\dot\mom=\advcoef\partial_\yco\mom+\wspeed^2\partial_\yco^2\field$.
Eliminating $\mom=\dot\field-\advcoef\partial_\yco\field$,
\begin{equation}
    \partial_\Ny^2\field - 2\advcoef\,\partial_\Ny\partial_\yco\field + (\advcoef^2-\wspeed^2)\,\partial_\yco^2\field = 0 .
    \label{eq:hyperbolic-reduction}
\end{equation}
Its principal symbol is the quadratic form $M^{ab}\xi_a\xi_b$ with
\begin{equation}
    M = \begin{pmatrix} 1 & -\advcoef \\ -\advcoef & \advcoef^2-\wspeed^2 \end{pmatrix} ,
    \qquad
    \det M = -\wspeed^2 < 0 ,
    \label{eq:optical-metric}
\end{equation}
an inverse optical metric of Lorentzian signature: the operator is
hyperbolic wherever $\wspeed>0$, for either sign of $\advcoef$. In this
form $\advcoef$ is a shift vector: unit lapse, and the advection
entirely in the shift. Characteristics are the null curves of $M$. The
null covectors are $\xi_a=(\advcoef\pm\wspeed,\,1)$; raising the index with
$M$ gives the tangent directions, because a null covector is orthogonal
to itself,
\begin{equation}
    \frac{\d\yco}{\d\Ny} = -\advcoef \pm \wspeed .
    \label{eq:characteristic-speeds}
\end{equation}
Prefactors multiplying the whole equation, such as $\gpref$, move affine
parameters along the characteristics but never the cone.

At the core, $\Deltas_{\rm loc}(1,\Ny)=\Deltas(\Ny)$ exactly
(Section~\ref{sec:gradient-operator}), so $\wspeed(+1)=2/\Deltas$ with no
approximation, and $\advcoef(+1)=2\dot{\Deltas}/\Deltas$. Hence
\begin{equation}
    \frac{\advcoef(+1)}{\wspeed(+1)} = \dot{\Deltas} ,
    \qquad
    \frac{\d\yco}{\d\Ny}\bigg|_{\yco=+1} = \frac{2}{\Deltas}\big(-\dot{\Deltas}\pm1\big) ,
    \label{eq:core-speeds}
\end{equation}
which on the homogeneous rate $\dot{\Deltas}=1-\epscore$ are
$2\epscore/\Deltas$ (outgoing, $\d\yco/\d\Ny>0$) and
$-2(2-\epscore)/\Deltas$ (incoming). A characteristic is incoming at
$\yco=+1$ exactly when its $\yco$-velocity is negative. The count is
therefore one incoming characteristic for $|\dot{\Deltas}|<1$, that is
$0<\epscore<2$ at the homogeneous level; two for $\dot{\Deltas}>1$ and
none for $\dot{\Deltas}<-1$. The window is unconditional for physical
trajectories, and this is a consequence of the core anchor
(Section~\ref{sec:coordinate}). One incoming characteristic admits exactly
one boundary condition.

Two features of \eqref{eq:core-speeds} matter. The outgoing speed
$2\epscore/\Deltas$ vanishes as $\epscore\to0$: in slow roll the core
boundary is almost exactly characteristic from the interior side, which
is physical (the horizon of de Sitter is marginally trapped) and is the
reason the core closure is delicate. Both speeds scale as $1/\Deltas$, so
their ratio $2/\epscore$ is independent of $\alpha$; the near-tangency
cannot be regularized by $\alpha$. At the outer edge $\advcoef(-1)=0$, the
cone is symmetric, $\d\yco/\d\Ny=\pm\wspeed(-1)=\pm(2/\Deltas)\e{-\Deltas_{\rm loc}(-1,\Ny)}$,
exponentially small at late times: one incoming, one outgoing, and the
outer edge barely communicates with the interior.

The elliptic argument of the earlier draft (a Cauchy problem for $\Lop$ at
fixed $\Ny$ is ill-posed, so Dirichlet at one end and Neumann at the other
is the well-posed configuration) is consistent with this: it governs the
spatial boundary value problem at fixed $\Ny$, whereas the hyperbolic
count governs the evolution in $(\yco,\Ny)$ and says how many conditions
the moving core admits.

\subsection{No data enters at the core}
\label{sec:no-data}

Write the first-order principal system in the variables
$(\partial_\yco\field,\mom)$. Its characteristic (Riemann) variables are
the left eigenvectors of the coefficient matrix,
\begin{equation}
    w_\pm \equiv \mom \pm \wspeed\,\partial_\yco\field ,
    \label{eq:riemann-variables}
\end{equation}
with $w_+$ carried along $\d\yco/\d\Ny=-\advcoef-\wspeed$ (incoming at
the core) and $w_-$ along $-\advcoef+\wspeed$ (outgoing). The energy flux
through the core in the norm of Section~\ref{sec:collocation-basis} is
\begin{equation}
    \frac{\volel}{4}\Big[(\wspeed+\advcoef)\,w_+^2 + (\advcoef-\wspeed)\,w_-^2\Big]_{\yco=+1} ,
    \label{eq:core-flux}
\end{equation}
so the incoming invariant enters with the positive coefficient
$\wspeed+\advcoef=(2/\Deltas)(1+\dot{\Deltas})\approx(2/\Deltas)(2-\epscore)$,
and the outgoing one with the negative coefficient $\advcoef-\wspeed$.

The Neumann condition \eqref{eq:bc-core} sets $w_+=w_-$: it is a
reflecting condition. It specifies the incoming invariant only in terms of
the outgoing one and supplies no data about the excised interior. This is
the correct closure. In inflation, perturbations are seeded at the horizon
by the noise; there is no inflow of information from the sub-horizon
interior that the model needs to capture. $\mom(1,\Ny)$ is not boundary
data: it is slaved to the interior by reflection, $\mom_\core=w_-$. Under
reflection the flux \eqref{eq:core-flux} collapses to
$\tfrac12\volel\advcoef(1)\mom_\core^2>0$, the work done by the moving
boundary as the comoving horizon shrinks past the fields. This is bounded
growth at the rate $\advcoef$, physical and not something to cancel. The
continuum problem is well posed with this growth; it is not energy
decaying at the core.

A second condition on $\mom_\core$ cannot be added. A data closure
$w_+=g$ and the reflecting closure $w_+=w_-$ compete for the same single
characteristic slot; imposing both over-determines the problem. On
10 July 2026 a fixed target for $\mom_\core$ was read as a model closure
supplying incoming data (the assertion that the excised interior looks
like the homogeneous instanton core). That reading was reversed on
11 July: no data enters at the core, and a penalty on $\mom_\core$ exists
only to cancel the piston flux, which a flat-norm strict-decay estimate
misreads as an instability. The over-determination was the origin of a
spurious stiffness growing as the square of the number of collocation
points (Section~\ref{sec:collocation-basis}).

The characteristic variable on which the single boundary penalty acts is
the incoming invariant
\begin{equation}
    \win = \mom_\core + \frac{2}{\Deltas}\,(\partial_\yco\field)_\core ,
    \label{eq:w-in}
\end{equation}
with penalty strength set by the flux coefficient
$(\wspeed+\advcoef)(1)=(2/\Deltas)(1+\dot{\Deltas})$. The relative
weighting of the two terms in $\win$ is set by $\Deltas$ alone, with no
dependence on $\epscore$: at early times, where $\Deltas$ is of order
$\ln(1+\alpha)$, the invariant is dominated by the gradient term, and at
late times, where $\Deltas\sim10$, by $\mom_\core$. A penalty on
$(\partial_\yco\field)_\core$ alone or on $\mom_\core$ alone is the
characteristic projection in neither limit.

\subsection{The response sector}

The response equations \eqref{eq:inst-rphi}--\eqref{eq:inst-rpi} have the
same principal symbol: $\Lop$ acts on $\rmom$ in the $\dot{\rfield}$
equation, the mirror image of the forward sector, and eliminating
$\rfield$ gives \eqref{eq:hyperbolic-reduction} for $\rmom$ with the same
$\advcoef$ and $\wspeed$. The adjoint of a hyperbolic operator has the
same characteristics, traversed in the opposite direction because the
response sector is integrated backward in $\Ny$. In backward time the
velocities reverse, so the roles swap: the mode that was incoming for the
forward sector is outgoing for the response sector and conversely. Again
exactly one characteristic enters at the core, and the same window
$0<\epscore<2$ governs both sectors. The response sector's incoming mode
is the slow, nearly characteristic one, so the delicacy of the closure is
if anything greater on this side.

The response condition is homogeneous. The boundary term of the action
at the core is affine in $(\field,\mom)$ with coefficients linear in
$(\rfield,\rmom)$; its variation with respect to $(\field,\mom)$ carries
no target. The response fields are Lagrange multipliers, their only
inhomogeneous datum is the terminal source, and there is no
``what lies inside the horizon'' question on this side. The explicit form
follows from the self-adjointness boundary term of
Section~\ref{sec:msr-action} with the operator acting on $\gpref\rmom$:
the natural condition is $\partial_\yco(\gpref\rmom)=0$ at $\yco=+1$, which
reduces to $\partial_\yco\rmom=0$ when $\gpref$ is $\yco$-independent at
the core. Since $\partial_\yco\mom$ need not vanish at the core
($\mom=\dot\field-\advcoef\partial_\yco\field$ gives
$\partial_\yco\mom|_\core=-\advcoef_\core\partial_\yco^2\field|_\core$),
$\gpref'$ is generically nonzero there and the two differ. In the scheme
of Section~\ref{sec:collocation-scheme} the response boundary rows are
generated from the discrete Hamiltonian and need not be written down; the
continuum statement is recorded for the indicial analysis and for the
paper.

\subsection{Temporal boundary conditions}

\begin{align}
    \field(\yco,\Ninit) & = \phiinit ,
    &
    \mom(\yco,\Ninit) & = \mom_{\mathrm{SR,init}} ,
    &
    \text{(all } \yco \text{)}
    \label{eq:bc-init}
    \\
    \field(1,\Nfinal) & = \phiend ,
    & & &
    \label{eq:bc-final-core}
    \\
    \rfield(\yco,\Nfinal) & = 0 ,
    &
    \rmom(\yco,\Nfinal) & = 0 ,
    &
    \text{(}\yco<1\text{, free)}
    \label{eq:bc-final-free}
\end{align}
with the constraint \eqref{eq:bc-final-core} enforced via a Lagrange
multiplier $\lag$ that sources $\rfield(\yco,\Nfinal)$ as a point
contribution at $\yco=1$; see Section~\ref{sec:collocation-scheme} for
the resulting terminal condition.
No boundary condition is imposed at $\yco<1$ forcing $\field$ to reach
$\phiend$: we do \emph{not} impose an explicit `freezing' of
$\field(\yco,\Ny)$ once it crosses $\phiend$ (see
Section~\ref{sec:zeta-extraction}).

\begin{panel}
\semibold{Linearity in $\lag$ is a tree-level property.} When the
response equations are linear and homogeneous in $(\rfield,\rmom)$, the
terminal source is the only inhomogeneity and the response solution scales
exactly as $\rfield=\lag\hat{\rfield}$ with $\hat{\rfield}$ independent of
$\lag$. This is guaranteed by the MSR construction at tree level, that
is, for Gaussian noise and the saddle-point approximation. It fails in two
ways: when the noise amplitude depends on the state
($\partial\Sigvol{}/\partial u\neq0$, the Riccati term of
Section~\ref{sec:instanton-eqs}), and when loop corrections from a
Schwinger--Keldysh effective action are included. Any use of the
rescaling, for instance to control the dynamic range of the response
fields, is contingent on linearity and should be labelled as such.
\end{panel}
\section{No explicit freezing; extraction of $\zprofile(\yco)$}
\label{sec:zeta-extraction}
% ======================================================================

In the single-trajectory instanton, `end of inflation' is a single event.
Here, different $\yco$ reach $\phiend$ at different times
$\Ny_{\mathrm{stop}}(\yco)$, with $\Ny_{\mathrm{stop}}(-1) = $ (the noiseless
crossing time) and $\Ny_{\mathrm{stop}}(1) = \Nfinal$. Treating
$\Ny_{\mathrm{stop}}(\yco)$ as an a priori stopping boundary makes this a
free-boundary problem: the domain of \eqref{eq:instanton-eqs} would have a
curved upper edge that is itself part of the solution.

We avoid this by \emph{not} imposing any stopping or freezing condition:
the system \eqref{eq:instanton-eqs} is integrated over the full rectangle
$\yco\in[-1,1]$, $\Ny\in[\Ninit,\Nfinal]$, with the single core condition
\eqref{eq:bc-final-core}, and $\zprofile(\yco)$ is read off as a
\emph{post-processing} step applied to the resulting $\field(\yco,\Nfinal),
\mom(\yco,\Nfinal)$.

\par\vspace{0.5em}\noindent
\semibold{Scale assignment: density matching, not a crossing scan.}
The instanton can, in principle, leave each shell with its own
isocurvature perturbation relative to the others (small and rapidly
decaying in a single-field model, but not exactly zero), so the physically
correct radius assignment compares surfaces of \emph{fixed energy density}
between the shell and the background, not surfaces of fixed field value on
a fixed-$\Ny$ (flat-gauge-like) slicing -- these differ once $\zprofile$ is
$\Or(1)$, in the same way uniform-density and flat time-slicings differ in
the standard $\delta N$ formalism. Concretely, for each $\yco$:
\begin{enumerate}
    \item Take the reconstructed state $\big(\field(\yco,\Nfinal),\mom(\yco,\Nfinal)\big)$
    at the \emph{shared} final time $\Nfinal$ -- the same $\Ny$ for every
    $\yco$, read directly off the global solution, with no per-shell
    crossing search.
    \item Continue \emph{that shell's own} trajectory noiselessly forward
    from this state until $\epsSR=1$ (genuine end of inflation), giving an
    endpoint $e$-fold $\Ny_{\mathrm{end}}(\yco)$ and an energy density
    $\rho_{\mathrm{end}}(\yco) \equiv \rho\big[\field_{\downarrow}(\yco,\Ny_{\mathrm{end}}),
    \mom_{\downarrow}(\yco,\Ny_{\mathrm{end}})\big]$ there, where the
    $\downarrow$ subscript denotes this per-shell noiseless downflow.
    \item Density-match against the background: find $\Ny_{\mathrm{nl}}(\rho_{\mathrm{end}}(\yco))$,
    the $e$-fold at which the noiseless trajectory itself has energy
    density $\rho_{\mathrm{end}}(\yco)$.
    \item
    \begin{equation}
        \zprofile(\yco) = \Ny_{\mathrm{end}}(\yco) - \Ny_{\mathrm{nl}}\big(\rho_{\mathrm{end}}(\yco)\big) .
        \label{eq:zeta-extraction}
    \end{equation}
\end{enumerate}
This refines the construction used for the single-trajectory case, which
matches against the background without the downflow to $\epsSR=1$; the
difference is negligible in single-field models and matters only once
isocurvature is present.
Applying it pointwise in $\yco$ here is the natural generalization,
sidesteps needing to surface-find $\Ny_{\mathrm{stop}}(\yco)$ at all
(addressing the free-boundary concern above by a different route than
originally intended: the downflow does the work that the crossing scan
was trying to do, without needing $\field(\yco,\Nfinal)$ to already be at
or past $\phiend$), and is expected to give monotonically ordered radii
across shells provided $\zprofile$ is not increasing so rapidly with
$\yco$ that it produces a Type~II perturbation.


% ======================================================================
\section{Physical and comoving scale assignment}
\label{sec:scale-assignment}
% ======================================================================

The construction above gives $\zprofile(\yco)$, a genuine profile rather
than the single number the existing single-trajectory pipeline produces.
Turning this into radius and mass estimates for
\texttt{CompactionFunction} (or the equivalent
container for \texttt{GradientCoupledInstanton})
requires distinguishing three different
notions of ``scale'' that the single-trajectory case did not need to
separate, because there a single comoving radius played all three roles
at once.

\subsection{Comoving radius}

$\rad(\yco,\Ny)$ is already known for every $\yco$, directly from the
coordinate map \eqref{eq:y-definition} evaluated at $\Ny=\Nfinal$ once
$\Deltas(\Nfinal)$ has been solved for --- no separate calculation. Being
comoving, it carries none of the metric perturbation: it is a coordinate
label, unperturbed by construction, in the same sense any comoving
coordinate is.

\subsection{Areal radius}

The GR-meaningful radius, and the one entering the standard compaction
function formula (Yoo 2022),
\begin{equation}
    C(\rad) = \frac23\Big[1-\big(1+\rad\,\zprofile'(\rad)\big)^2\Big] ,
    \label{eq:compaction-yoo}
\end{equation}
is the areal radius $R(\rad) = a\,\rad\,e^{\zprofile(\rad)}$: this formula
is derived directly from $\d R/\d\rad = a\,e^{\zprofile}(1+\rad\zprofile')$,
so the areal-radius relation is already built into it. Note
\eqref{eq:compaction-yoo} takes $\zprofile(\rad)$ and
$\zprofile'(\rad)\equiv\d\zprofile/\d\rad$ --- comoving derivatives --- as
input, \emph{not} $\zprofile(R)$: using the areal coordinate here would be
a genuine mismatch, not an equivalent relabelling. Since
$\zprofile(\yco_j)$ is already known at every collocation node and
$\rad(\yco_j,\Nfinal)$ from the coordinate map, $\zprofile'(\rad)$ follows
by the chain rule, $\zprofile'(\rad) = (\d\zprofile/\d\yco)/(\d\rad/\d\yco)$,
using the collocation differentiation matrix for the numerator and the
analytic $\d\rad/\d\yco=-\tfrac12\Deltas(\Nfinal)\rad(\yco,\Nfinal)$ for the
denominator --- no new derivative machinery.
Alternatively, since we will extract $\zeta(r)$ anyway, the
deriative can be computed directly.

$R(\rad)$ is not itself an observable; its role is as the correctly
GR-translated input against which a numerical-relativity-calibrated
threshold $C_{\rm th}$ can be compared, since NR collapse simulations
(Musco, Escriv\'a, Young, and collaborators) are set up and calibrated
directly in areal-radius/compaction-function language. As $\rad$ shrinks
towards the interior, if $\zprofile(\rad)$ grows steeply enough $R(\rad)$
can stop being monotonic in $\rad$ --- the Type~II perturbation regime
already flagged (Section~\ref{sec:zeta-extraction}), corresponding to
$C_{\max}\to2/3$, and the reason the existing pipeline computes $r_{\max}$
and $r_{\rm peak}$ from $C(\rad)$ alone rather than from any
$\bar C$-based estimate sensitive to the $e^{3\zprofile}$ volume
weighting.

\subsection{Physical (present-day) scale}

A third, independent notion: the wavelength a given comoving mode
corresponds to today, needed for the mass/scale axis of the final PBH
mass function. This uses the existing Leach--Liddle scale-matching
equation (astro-ph/0305263), solved \emph{once}, anchored at the fixed
outer edge $\rout$ --- exactly as for the single-trajectory case, since
$\yco=-1$ sits exactly on the noiseless background by construction
(\eqref{eq:bc-outer}), giving a physical scale $\rad_{\mathrm{phys,out}}$. Since comoving
wavenumber satisfies $k\propto1/\rad$, every other shell's physical scale
follows by a fixed ratio,
\begin{equation}
    \rphys(\yco) = \frac{\rad(\yco,\Nfinal)}{\rout}\;\rad_{\mathrm{phys,out}} ,
    \label{eq:rphys-ratio}
\end{equation}
with no further per-shell Leach--Liddle solve and no per-shell inversion
of the background expansion history: comoving-ness already fixes the
today's-scale conversion to a single universal constant, common to every
mode regardless of when it crossed the horizon, and that constant is
exactly what solving Leach--Liddle once determines. The genuine
per-shell/isocurvature-sensitive calculation --- the downflow from each
shell's own instanton endpoint to $\epsSR=1$, giving that shell's own
$V_{\rm end,\,downflow}$ and hence its own duration in the sense already
used for $\zprofile(\yco)$ (Section~\ref{sec:zeta-extraction}) --- is
unaffected and answers a different physical question (how long \emph{that
particular shell} takes to finish inflating) from "which comoving scale is
this" (a purely geometric fact, answered once and for all by
\eqref{eq:rphys-ratio}).

\begin{panel}
\semibold{A candidate refinement, considered and rejected.} One might
instead try to assign each shell its own horizon-crossing $e$-fold by
inverting $a(\Ny')\Hub(\Ny')=k[\rad(\yco,\Nfinal)]$ against the
\emph{background} trajectory. This was considered and is unnecessary: it
answers a question (``which $e$-fold did this comoving scale cross the
horizon'') that plays no role once $\rad$ is already comoving, since the
comoving-to-today conversion is a single constant for every mode by
definition, not something that depends on horizon-crossing time
individually. The single global ratio \eqref{eq:rphys-ratio} already
contains everything the inversion would have supplied.
\end{panel}

\subsection{Anchoring the physical scale}
\label{sec:scale-anchor}

The anchor is evaluated at the start of the transition, $\Ninit$, not by a
downflow from the endpoint. Since $\yco=-1$ is pinned to the noiseless
background throughout, its evolution after $\Ninit$ is the remaining
noiseless trajectory and needs no integration; the number of $e$-folds
from the start of the transition to the true end of inflation is $\Ninit$
itself. The matching equation is written in terms of the Hubble rate at
horizon crossing and at the end of inflation, $\Hub_k$ and $\Hub_{\rm end}$,
rather than in terms of $\Vpot$ and $\epsSR$: the Friedmann relation
connecting them is the potential's responsibility and should not be
assumed to take a particular form. Because $\rout$ is $(1+\alpha)$ larger
than the horizon at $\Ninit$, the anchored scale carries the factor
$(1+\alpha)$; without it the assigned scale would drift with $\alpha$.
An earlier form of the matching equation, which merged the
horizon-crossing Hubble rate with the end-of-inflation energy density in a
single logarithm, predicted a physical scale varying as $\Hub^{1/2}$ at
fixed $\Ny$, contradicting $k=a\Hub$; it was corrected on 6 July 2026,
and the single-trajectory construction was re-anchored at the outer edge
and worked forward, with the $e$-fold count before the end of inflation
obtained by arithmetic from $\Ninit$. The equivalence question of the
earlier draft (whether the discrete peeling construction uses the
instanton's own local trajectory for the core's horizon) was answered in
the affirmative: it does, and the two schemes agree at the core.

\subsection{The core scale, and a difference between pipelines}
\label{sec:core-scale}

Under the core anchor the comoving ratio of \eqref{eq:rphys-ratio} at the
core is $\rad(1,\Nfinal)/\rout=\e{-\Deltas(\Nfinal)}$ with the core's own
$\Deltas$, and the physical scale assigned to the core is that of the
core's own horizon at $\Nfinal$; the factor $(1+\alpha)$ cancels between
$\rout$ and the anchor. The single-trajectory pipeline assigns the core a
scale built from the perturbed endpoint, through both the $e$-fold count
(which includes $\dNstar$) and the perturbed Hubble rate at crossing. The
two schemes make different choices about what $\delta N$ does to the
scale: whether the excess $e$-folds relabel the scale or appear as an
amplitude at a fixed scale. Both are internally consistent
separate-universe bookkeepings; they coincide only as $\dNstar\to0$. The
difference is one of mass attribution, not of threshold: the collapse
criterion is amplitude-driven and insulated, but any statement of the form
``this fluctuation makes a black hole of mass $M$'' differs between the
two by a $\dNstar$-dependent shift in $\ln k$. The intermediate-scale
discrepancy between the pipelines (of order $\e{\dNstar}/(1+\alpha)$ at
the outer layer) was attributed on 6 July 2026 to the separate-universe
approximation being on the verge of breaking down, and is unquantified.

\begin{todo}
Which scale assignment to adopt for the core, and how to reconcile the
mass attribution between the gradient-corrected and single-trajectory
pipelines, is open. The comparison recorded on 14 July 2026 was made under
the background anchor and should be redone under the core anchor. See
\texttt{RECONSTRUCTION.md}, Open~6.
\end{todo}
\section{Collocation basis: Legendre--Gauss--Lobatto nodes}
\label{sec:collocation-basis}
% ======================================================================

Rather than expand in eigenfunctions of $\Lop$ (not available in closed
form once $\Lop$ carries the drift term of Section~\ref{sec:gradient-operator}),
we
represent $\field,\mom,\rfield,\rmom$ by their values on a fixed grid of
collocation points and work directly with differentiation matrices. This
was chosen over Chebyshev collocation specifically so that the
self-adjoint structure of $\Lop$ (Section~\ref{sec:msr-action}) can be
preserved numerically via a summation-by-parts (SBP) operator with a
\emph{diagonal} norm matrix --- the standard construction for
Legendre--Gauss--Lobatto (LGL) nodes, not available for Chebyshev without
retrofitting a non-diagonal norm.

\subsection{Nodes, weights, and differentiation matrices}

The $\modeidx_{\max}+1$ LGL nodes $\{\yco_j\}$ on $[-1,1]$ are $\pm1$
together with the interior roots of $P_{\modeidx_{\max}}'(x)=0$ (the
derivative of the Legendre polynomial of degree $\modeidx_{\max}$),
computed via the eigenvalues of a symmetric tridiagonal Jacobi matrix ---
standard, tabulated, and readily available (e.g.\ the Golub--Welsch-type
routines in \texttt{FastGaussQuadrature.jl}, or the reference Python
implementation at \texttt{sphglltools}); no closed form exists, unlike the
Chebyshev $\cos(j\pi/\modeidx_{\max})$ nodes, but this is a solved problem
requiring no new derivation. The associated LGL quadrature weights are
\begin{equation}
    w_j = \frac{2}{\modeidx_{\max}(\modeidx_{\max}+1)\,\big[P_{\modeidx_{\max}}(\yco_j)\big]^2} ,
    \label{eq:lgl-weights}
\end{equation}
with $w_j\cdot\measure(\yco_j,\Ny)$ the discrete quadrature weight for any
integral against the physical measure $\measure(\yco,\Ny)\,\d\yco$ of
Section~\ref{sec:msr-action} --- in particular for the endpoint,
$w_{\modeidx_{\max}} = 2/[\modeidx_{\max}(\modeidx_{\max}+1)]$, using
$P_{\modeidx_{\max}}(1)=1$: a known, finite number, needed below for the
terminal condition. Standard LGL differentiation matrices $D,D^{(2)}$
(first and second derivative on the nodes) are built once from the grid
alone --- they do not depend on $\Deltas(\Ny)$ or any other physics input,
so are computed once at the start and reused at every $\Ny$.

\subsection{Discretizing $\Lop$}

$\Lop$ (\eqref{eq:Lop-definition}) has $\Ny$-dependent, but not
$\Ny$-dependent-\emph{basis}, coefficients: at each $\Ny$, discretizing it
is
\begin{equation}
    \big(\Lop\field\big)_j
    =
    \e{\Deltas(\Ny)}\e{\Deltas(\Ny)\yco_j}
    \left[\frac{4}{\Deltas(\Ny)^2}\big(D^{(2)}\field\big)_j - \frac{2}{\Deltas(\Ny)}\big(D\field\big)_j\right] ,
\end{equation}
matching \eqref{eq:Lop-definition} --- note $\Lop$ itself does \emph{not}
carry the $1/\rout^2$ prefactor; that is applied separately, together with
$1/\aHloc^2$, wherever $\Lop$ appears in the equations of motion (e.g.\
\eqref{eq:inst-pi}). This separately-applied prefactor is evaluated node-by-node as
$\aHloc(\yco_j,\Ny)$ throughout the collocation scheme, i.e.\ it is already
the local-$\Deltas$ object of the panel in Section~\ref{sec:coordinate}
($1/[\rout^2\aHloc(\yco_j,\Ny)^2]=\e{-2\Deltas_{\rm loc}(\yco_j,\Ny)}$); no
change to the numerics follows from that remark, it is purely a
bookkeeping identification. The $\Deltas(\Ny)$ appearing \emph{inside}
$\Lop$ above, by contrast, is the single core-based value shared by every
node.
i.e.\ $D,D^{(2)}$ applied to the node-value vector, then multiplied by
node-dependent \emph{scalar} coefficients recomputed at each $\Ny$ --- the
same pattern already used for every other solution-dependent coefficient
in this scheme ($\Hub^2\atloc$, $\epsSR\atloc$, $\Dcm{ij}\atloc$), no new
numerical primitive. The advection term $\advcoef(\yco,\Ny)\partial_\yco(\cdot)$
is likewise just $\mathrm{diag}(\advcoef(\yco_j,\Ny))\,D$ applied to the
node-value vector.

\subsection{The energy method and the choice of norm}
\label{sec:energy-method}

A first-derivative operator $D$ on the grid is summation-by-parts (SBP)
with respect to a symmetric positive-definite norm matrix $\Hnorm$ if
\begin{equation}
    \Hnorm D + (\Hnorm D)^{\transpose} = B , \qquad B = \mathrm{diag}(-1,0,\dots,0,+1) ,
    \label{eq:sbp-identity}
\end{equation}
the discrete mirror of integration by parts. In continuum language
$\Hnorm$ is a kernel whose composition with $D$ is a total derivative.
This is a property of $D$; it has nothing directly to do with physics,
and the inner product $\langle v,u\rangle=v^{\transpose}\Hnorm u$ need not
have a physical meaning. We call $\Hen=u^{\transpose}\Hnorm u$ the
$\Hnorm$-energy. It is a control functional, not a physical energy, and it
is not $\Hfp$. For the LGL nodes with $\Hnorm=\mathrm{diag}(w_j)$ the
identity \eqref{eq:sbp-identity} holds to machine precision.

The method has three steps. First, choose $\Hnorm$ so that the interior
contribution to $\d\Hen/\d\Ny$ is sign-definite: zero for a first-order
skew operator (the total-derivative case), or $-u^{\transpose}Mu\leq0$ with
$M\succeq0$ for the second-order dissipative piece. For a generic operator
no such $\Hnorm$ need exist, and for a given operator it is essentially
unique. Second, read the remaining boundary flux in characteristic
variables; it splits into pieces of definite sign attached to outgoing
characteristics, which drain $\Hen$ and must not be touched, and to
incoming characteristics, which must be controlled. Third, supply exactly
the incoming data. Then $\d\Hen/\d\Ny\leq2\beta\Hen$ for a fixed rate
$\beta$, and $\Hen(\Ny)\leq\e{2\beta\Ny}\Hen(0)$. Well-posedness of the
continuum problem is the existence of such an $\Hnorm$ and a set of
boundary conditions; the method engineers $\Hen$ rather than discovering
a conserved quantity. Too few conditions and the incoming flux is
uncontrolled; too many and the problem is over-determined. The discrete
scheme mirrors this: SBP reproduces the first step exactly, and a
simultaneous approximation term (SAT) reproduces the second and third by
supplying the incoming data weakly, with a penalty strength chosen so
that the discrete boundary flux equals the continuum one.

Positive-definiteness of $\Hnorm$ is what makes the method say anything
about $u$. It gives norm equivalence, $c\|u\|^2\leq\Hen\leq C\|u\|^2$, so a
bound on $\Hen$ is a bound on $\|u\|$, and any two admissible norms are
mutually equivalent, so they agree on whether $u$ is bounded even though
their rates differ. There is no paradox of an $\Hen_1$-bounded,
$\Hen_2$-unbounded solution. The constant $C/c$ is the condition number of
$\Hnorm$; for the weighted norm below it is large near $\Ninit$, where the
measure is sharply graded, so the bound is loose precisely in the stiff
stretch.

For the onion operator only the weighted norm
$\Hnorm_\volel=\mathrm{diag}\big(w_j\volel(\yco_j,\Ny)\big)$ makes the
interior sign-definite, because $\Lop$ is self-adjoint with respect to
$\volel\,\d\yco$ and not with respect to the flat measure. In the flat
norm the interior contribution to $\d\Hen/\d\Ny$ does not vanish, the
leftover source sits on the core node, and $\Hen$ is not a usable control
functional at all. The $\volel$ in the characteristic flux
\eqref{eq:core-flux} is the same $\volel$. The measure is physical, being
handed down by the action; the Lyapunov functional built from it is not.

$\Hfp$ is indefinite and can never be a control functional. It is,
however, a bounded quadratic form in the joint state
$u=(\field,\mom,\rfield,\rmom)$, so $|\Hfp|\leq(M/c)\,\Hen$ once $\Hen$ is
bounded on that joint state. Boundedness of $\Hfp$ therefore follows by
domination from an energy estimate on both sectors together, with no
reference to the flux structure of $\Hfp$ itself, and is immune to any
truncation of the response sector. That the same incoming data regularize
both $\Hen$ and $\Hfp$ is because both time derivatives are reduced by
integration by parts against the same operator, so both throw off
boundary forms in the same characteristic basis. The estimate requires an
energy on the response sector as well as the forward one; a response
sector closed by hard elimination has none.

\begin{panel}
\semibold{Norm for stability, transpose for adjointness, never mixed.}
Two different calculations involve the boundary. The energy estimate that
fixes the penalty strength lives in $\Hnorm_\volel$. The response equations
are the variation of the same discrete action with respect to the field
values, in which the transpose of the forward Jacobian appears
automatically because the response fields are covectors; no norm enters.
An earlier flat-norm derivation of the penalty, and an earlier attempt to
obtain the response closure by applying the forward recipe with a sign
flip, each mixed the two.
\end{panel}

\subsection{Boundary conditions}
\label{sec:colloc-bcs}

At $\yco=-1$ ($j=0$): ordinary Dirichlet rows for $\field,\mom$ (fixed to
$\phinl(\Ny),\pinl(\Ny)$) and $\rfield,\rmom$ (fixed to $0$). Since
$\advcoef(-1)=0$ exactly, this edge carries no destabilizing energy term,
and strong imposition is left unchanged.

At $\yco=+1$ ($j=\modeidx_{\max}$): the target scheme imposes the single
admissible forward condition weakly, as one SAT acting on the incoming
Riemann invariant \eqref{eq:w-in},
\begin{equation}
    \frac{\d}{\d\Ny}(\cdot)_{\modeidx_{\max}}
    \;\mathrel{+}=\;
    -\frac{\tau}{w_{\modeidx_{\max}}\volel(1,\Ny)}\,\ell\,\big(\win-\win^{\rm target}\big) ,
    \label{eq:single-sat}
\end{equation}
with $\ell$ the left eigenvector of the incoming characteristic in the
$(\partial_\yco\field,\mom)$ variables, $\tau$ fixed by requiring that
the discrete boundary flux in $\Hnorm_\volel$ equal the continuum flux
\eqref{eq:core-flux}, and the target the reflecting value $w_-$ (no data).
The penalty is derived as a boundary term of the discrete action
(Section~\ref{sec:collocation-scheme}), so that the response boundary rows
follow by differentiation and the two sectors are discretely adjoint. The
core node $\field_{\modeidx_{\max}}$ is a free, integrated degree of
freedom; $\mom(1,\Ny)$ and $\rfield(1,\Ny)$ carry no boundary row and are
determined by the dynamics and the shooting problem.

\begin{todo}
Nothing in this subsection is implemented. The scheme in production as of
October 2026 closes the forward sector with two penalties at the core (one
on $\field_{\modeidx_{\max}}$ toward a live regularity target, one on
$\mom_{\modeidx_{\max}}$ toward a fixed target), both derived in the flat
LGL norm, and closes the response sector by hard elimination of
$\partial_\yco\rmom=0$. The penalty on $\mom$ was introduced to cancel the
piston flux $\tfrac12\volel\advcoef\mom_\core^2$, which a flat-norm
strict-decay estimate reads as an instability; the live target on $\field$
is, for LGL nodes, a derivative-type condition
($w_{\modeidx_{\max}}D_{\modeidx_{\max},\modeidx_{\max}}=\tfrac12$ exactly).
Two penalties on a one-characteristic boundary are an over-determination
whose strength grows as $\tau/w_{\modeidx_{\max}}\propto\modeidx_{\max}^2$,
and the fixed target does not vanish at convergence, so the converged
solutions depend materially on $\tau$. Hard elimination in the response
sector makes the two sectors non-adjoint. Replacing all of this with the
single characteristic penalty in $\Hnorm_\volel$ and a differentiated
response sector is Phase~3 of the plan; see \texttt{RECONSTRUCTION.md},
Part~C and Open~2.
\end{todo}

\begin{panel}
\semibold{Historical: hard elimination and the first SAT.} The first
implementation imposed the Neumann condition by hard elimination,
solving $\sum_k D_{\modeidx_{\max},k}\field_k=0$ for the boundary value and
substituting it into the interior rows. Elimination breaks the SBP
identity and forfeits the energy estimate; it produced a spurious
right-half-plane spectrum whose abscissa grew as
$\modeidx_{\max}^{1.6}$ at every $\Deltas$, independent of the integrator.
The remedy was split-form advection with an SAT at the core, which bounded
the abscissa. The derivation used the flat norm and two penalties, which is
the over-determination described above.
\end{panel}
% ======================================================================
\section{Collocation scheme}
\label{sec:collocation-scheme}
% ======================================================================

\subsection{Grid representation}

$\field,\mom,\rfield,\rmom$ are represented directly by their node values
$\field_j(\Ny)\equiv\field(\yco_j,\Ny)$, etc. --- no mode expansion, and no
lift against the noiseless trajectory is structurally required. The
Dirichlet boundary row simply sets $\field_0(\Ny)=\phinl(\Ny)$
(Section~\ref{sec:collocation-basis}). Tracking the deviation
$\field_j-\phinl$ rather than $\field_j$ may still be worth doing for
conditioning, but this is an optional choice and is deferred to the
companion planning document.

\subsection{The discrete Hamiltonian}
\label{sec:discrete-hamiltonian}

The equations of motion are not obtained by discretizing
\eqref{eq:instanton-eqs}. They are Hamilton's equations of a single
discrete Hamiltonian $\Hh$, constructed from $\Hfp$ as follows.
Integrate by parts in $\yco$ in the continuum first, so that the gradient
term contains only first derivatives,
\begin{equation}
    -\int\volel\,\d\yco\,\rmom\,\gpref\,\Lop\field
    = \Cpref\kappa\int\d\yco\,\weightfn\,\partial_\yco(\gpref\rmom)\,\partial_\yco\field
    - \Cpref\kappa\big[\weightfn\,\gpref\rmom\,\partial_\yco\field\big]_{-1}^{+1} ,
\end{equation}
with $\kappa$, $\weightfn$ as in \eqref{eq:hfp-flux}. The SBP identity
\eqref{eq:sbp-identity} mirrors integration by parts at first-derivative
level only; there is no equally clean identity for a second-derivative
matrix, so discretizing the first-order form is what makes the discrete
gradient operator symmetric by construction, with its boundary term
explicit. Then restrict every field to its node values, replace
$\int\volel\,\d\yco\to\sum_jw_j\volel(\yco_j,\Ny)$ and
$\partial_\yco\to D$, and evaluate every pointwise nonlinearity
($\Hub^2$, $\epsSR$, $\Vpot'$, $\Sigvol{ij}$, $\gpref$) at the node. The
result is a scalar function
\begin{equation}
    \Hh\big(\{u_j\},\{P_j\};\Ny\big) ,
    \qquad u_j=(\field_j,\mom_j) , \quad P_j = \big(w_j\volel_j\rfield_j,\; w_j\volel_j\rmom_j\big) ,
\end{equation}
with the canonical momenta carrying the quadrature weights, as the
continuum momenta carry $\volel$ \eqref{eq:hfp-canonical}. The quadrature
is only exact to a finite degree, so $\Hh$ differs from $\Hfp$ by an
aliasing error; that error is part of the variational object and must
not be removed by over-integration, which would break the compatibility
of $\Hnorm$ and $D$. The state dependence of $\Deltas$, $\dot{\Deltas}$,
$\advcoef$, $\volel$ and $\ncount$ through the core node is inside $\Hh$,
evaluated from $u$ at each call by \eqref{eq:Deltas-closed} and
\eqref{eq:Deltas-dot-closed}; it is not frozen.

The forward flow is $\dot u_j=\partial\Hh/\partial P_j$. It is linear in
$P$ and short: it is \eqref{eq:inst-phi}--\eqref{eq:inst-pi} at the
nodes, with $\partial_\yco\to D$, the gradient term in its symmetric
first-order form, and the SAT of \eqref{eq:single-sat} if the SAT is a
term in $\Hh$. The response flow is $\dot P_j=-\partial\Hh/\partial u_j$.
It is obtained by differentiating $\Hh$, not by transcription: every
derivative of Section~\ref{sec:instanton-eqs}, including the rank-one
couplings through the core node and the quotient rule in
\eqref{eq:Deltas-dot-closed}, is then a chain rule applied by a machine
rather than by hand, and the factor of $\tfrac12$ on the Riccati term is
automatic. Discrete adjoint consistency holds by construction, the
balance law \eqref{eq:balance-law} holds with $R=0$, and the Maupertuis
relation \eqref{eq:maupertuis} is exact for the discrete problem. The
interior rows of the differentiated response flow converge to the
discretized continuum equations \eqref{eq:inst-rphi}--\eqref{eq:inst-rpi}
plus the full-variation terms as the grid is refined; the boundary rows
differ at $\Or(1)$, and the discrete rows are the correct ones for the
discrete problem.

\para{Acceptance test} The Jacobian of the joint flow $z=(u,P)$ must be
Hamiltonian, $J_{\rm full}=\Omega\,\mathrm{Hess}\,\Hh$, equivalently
\begin{equation}
    \big(\Omega J_{\rm full}\big)^{\transpose} = \Omega J_{\rm full} ,
    \qquad
    J_{\rm full} = \begin{pmatrix} A & \Sigvol{} \\ C & -A^{\transpose} \end{pmatrix} ,
    \quad \Sigvol{}=\Sigvol{}^{\transpose} , \ C=C^{\transpose} ,
    \label{eq:hamiltonian-check}
\end{equation}
with $A=\partial b/\partial u+(\partial\Sigvol{}/\partial u)P$. This is
three block identities and should hold to machine precision for any
twice-differentiable $\Hh$. It replaces the earlier test
$J_{\rm resp}+J_{\rm fwd}^{\transpose}=0$, which is valid only when
$\partial\Sigvol{}/\partial u=0$.

\begin{todo}
Whether the SAT is a term in $\Hh$ (so that dual consistency holds by
construction and the response boundary rows follow from the same
differentiation) or is kept outside $\Hh$ and added to the assembled
right-hand side (so that $\Hh$ stays free of non-analytic operations such
as the absolute value in the penalty strength, and the Hamiltonian check
\eqref{eq:hamiltonian-check} is run on the unpenalized operator where it
holds exactly) is unresolved. A penalty with a strength fixed by a
frozen-coefficient energy argument in the wrong norm is in general not
variational; ``is the closure derivable from a boundary term in the
discrete action?'' is a design constraint on the single characteristic
penalty. See \texttt{RECONSTRUCTION.md}, Open~2.
\end{todo}

\subsection{Terminal condition on the boundary node}
\label{sec:terminal-collocation}

The Lagrange-multiplier construction of Section~\ref{sec:bcs} carries
over directly. Varying $S_{\mathrm{tot}}=S+\lag(\field_{\modeidx_{\max}}(\Nfinal)-\phiend)$
with respect to the node values $\{\field_j(\Nfinal)\}$ gives, using the
quadrature weights of \eqref{eq:lgl-weights},
\begin{equation}
    \sum_j w_j\,\volel(\yco_j,\Nfinal)\,\rfield_j(\Nfinal)\,\delta\field_j
    + \lag\,\delta\field_{\modeidx_{\max}}
    = 0
    \qquad \text{for all admissible } \{\delta\field_j\} .
\end{equation}
Since $\{\delta\field_j\}$ are independent, ordinary finite-dimensional
variables, this holds for every $j$ separately:
\begin{equation}
    \rfield_j(\Nfinal) = 0 \quad (j\neq\modeidx_{\max}) ,
    \qquad
    \rfield_{\modeidx_{\max}}(\Nfinal) = -\frac{\lag}{w_{\modeidx_{\max}}\,\volel(1,\Nfinal)} ,
    \label{eq:terminal-colloc}
\end{equation}
using $w_{\modeidx_{\max}}=2/[\modeidx_{\max}(\modeidx_{\max}+1)]$ from
\eqref{eq:lgl-weights}. Eq.~\eqref{eq:terminal-colloc} is the exact,
complete terminal condition. In the canonical variables of
Section~\ref{sec:hfp} it reads $P_{\field,\modeidx_{\max}}(\Nfinal)=-\lag$:
the terminal response momentum at the core is the multiplier itself. The same measure compensation appears in the
reduction to the gradient-free instanton: a response field supported at
the core must be written
\begin{equation}
    \rfield(\yco,\Ny)=\frac{\delta(\yco-1)\,\rfield_\core(\Ny)}{\volel(1,\Ny)}
    \label{eq:compensated-delta}
\end{equation}
for the pairing to reproduce the single-trajectory action.

\subsection{Shooting problem}
\label{sec:shooting}

Closed by a single scalar shooting condition: adjust $\lag$ so that
\begin{equation}
    \field_{\modeidx_{\max}}(\Nfinal) = \phiend .
\end{equation}
Structurally identical to the shooting problem already solved for the
single-trajectory instanton.

\subsection{Picard iteration}

\begin{enumerate}
    \item Forward pass: integrate $(\field_j,\mom_j)$ forward from
    \eqref{eq:bc-init} with $\rfield_j=\rmom_j=0$ (zeroth iterate).
    \item Backward pass: integrate $(\rfield_j,\rmom_j)$ backward from
    \eqref{eq:terminal-colloc}, using the current forward-pass solution to
    evaluate $-\partial\Hh/\partial u$.
    \item Forward pass: re-integrate $(\field_j,\mom_j)$ using the
    updated response-field values as sources.
    \item Repeat to convergence; adjust $\lag$ via an outer scalar
    root-find to satisfy the shooting condition.
\end{enumerate}
When the response flow is nonlinear in $P$ (Section~\ref{sec:instanton-eqs}),
the backward pass is a Riccati flow and the rescaling by $\lag$ is no
longer exact; the structure of the iteration is otherwise unchanged. A
global discrete variational solve (minimum-action methods) would subsume
the Picard alternation, but such methods struggle with degenerate noise,
and the massless de Sitter diffusion model is maximally degenerate
($\Dcm{12}=\Dcm{22}=0$); this is the reason for keeping the Picard
structure while making the two sectors discretely adjoint.
\section{Open issues}
\label{sec:open-issues}
% ======================================================================

\subsection{Sub-horizon noise: resolved}

An earlier draft carried this as the first, blocking open issue: treating
every $\yco$ interior to the current horizon as an independently
$\delta(\yco-\yco')$-correlated noise source manufactures degrees of
freedom that do not physically exist, since $\ncount(\yco,\Ny)\to0$ there
(fewer than one coarse-graining patch per shell). Three options were
weighed --- switching off noise there (fails the reduction-limit
cross-check), enforcing a single shared noise realization there (a hard,
time-varying shape constraint on the solution), and restricting $\yco$ to
the exterior of the horizon (previously rejected due to a
time-dependent eigenbasis) --- and a rank-1 shared-noise-channel
modification of the MSR action was proposed as a leading candidate,
explicitly flagged as unpropagated through the full equation set.

That proposal is now retired. Adopting the exterior-only, logarithmic
onion coordinate of Section~\ref{sec:coordinate}, combined with a
collocation (rather than explicit spectral) numerical method, removes the
objection to horizon-exclusion that motivated rejecting it originally,
and the problem is resolved \emph{by construction} rather than patched:
there is no sub-horizon region left in the domain to source spurious
independent fluctuations, so no special noise treatment is needed at all.
See Section~\ref{sec:coordinate} for the coordinate definition and the
panel explicitly narrating this reversal.

\subsection{Probability interpretation of $S_{\mathrm{MSR}}$ for the 2D configuration}

\begin{todo}
The single-trajectory instanton action $S_{\mathrm{MSR}}^{\mathrm{1D}}$
computes the probability of a specific event: the central Hubble patch
follows a given trajectory, \emph{with the surrounding shells assumed to
evolve noiselessly} (zero action cost by assumption). The two-dimensional
action $S_{\mathrm{MSR}}$ of \eqref{eq:msr-action} computes the probability
of the joint configuration $\field(\yco,\Ny)$ across the whole overdensity,
core and shells together. These are not simply refinements of the same
number: they are probabilities for two different events. Naively, one
might attempt to interpret the difference $S_{\mathrm{MSR}} -
S_{\mathrm{MSR}}^{\mathrm{1D}}$ as `the additional cost of the gradient
correction', but it was pointed out (correctly) in discussion that the
outer shells are not passive bystanders: they participate directly in the
compaction function and hence in the collapse criterion, exactly as they
do in the existing (noiseless-detachment) construction. So it is not
obviously right to think of the 1D action as already including a
`for-free' contribution from the outer layers that the 2D calculation
merely refines --- rather, the 1D action already implicitly describes the
joint configuration too, just with the outer-layer part of that
configuration fixed to be exactly noiseless. The correct comparison, and
the correct interpretation of what $S_{\mathrm{MSR}}$ computes once the
gradient coupling is switched on, has not yet been fully resolved and
needs further thought before results can be presented as `the corrected
PBH formation probability'.
\end{todo}

\subsection{The largest time equation for a stochastic (MSR) path integral}

\begin{todo}
The terminal condition $\rfield(\yco,\Nfinal)=0$ (Section~\ref{sec:bcs})
was justified heuristically by an analogue of the CTP `largest time
equation': placing the true response-field terminal condition at some
$\Ny_{\mathrm{turn}} > \Nfinal$ should not change the solution for
$\Ny \leq \Nfinal$, provided the interval $[\Nfinal,\Ny_{\mathrm{turn}}]$ is
genuinely noiseless and gradient-free, because the response-field
equations of motion are then linear and homogeneous there and, starting
from zero at $\Ny_{\mathrm{turn}}$, remain zero all the way back to
$\Nfinal$. This is physically reasonable and was checked at the level of
the linearized equations here, but the underlying general statement ---
that a Schwinger--Keldysh-type path integral, cut open at an intermediate
time with no later operator insertions, collapses to a functional delta
function enforcing agreement of the forward and backward fields at that
time --- does not appear to have a fully satisfactory published derivation
even in the standard quantum CTP case (see the parallel discussion in the
companion notes on the Dyson equation and the $\delta N$/Schwinger--Keldysh
correspondence), and has certainly not been established here for the
\emph{stochastic} (MSR) path integral in the presence of nonlinear noise
and gradient coupling. The quantum-case argument relies on unitarity of
the $[\Nfinal,\Ny_{\mathrm{turn}}]$ evolution; the analogous statement in
the stochastic case would rest on probability conservation, but the
precise functional-integral statement of this, and whether it implies the
same clean factorization, has not been derived carefully here. A
worthwhile check, cheaper than a full derivation: solve the existing
single-trajectory (1D) instanton with the terminal condition placed at
several $\Ny_{\mathrm{turn}} > \Nfinal$ and verify numerically that the
solution for $\Ny \leq \Nfinal$ is insensitive to the choice, before
relying on the same argument in the 2D case.
\end{todo}

% ======================================================================

\subsection{Items opened by the July 2026 analysis}

Each item below points to the corresponding entry of
\texttt{RECONSTRUCTION.md} (8 October 2026).

\begin{todo}
\semibold{Indicial analysis of the response sector at $\Deltas\to0$.} In
the $\volel$ convention the response sector has the same singular skeleton
as the forward sector (advection $1/\Deltas$, Laplacian $1/\Deltas^2$,
noise dilution $1/\Deltas$), which is the precondition for a Frobenius
analysis like \eqref{eq:frobenius-forward}. Posit
$\rfield=\Deltas^{p}\varrho_\field(\yco)+\dots$, find the indicial exponent,
and check that the leading singular terms cancel; the noise coupling
requires $p\geq1$ for the forward cancellation to be self-consistent. A
first attempt did not obviously close. Until it does, $\alpha$ cannot be
reduced on analytic grounds. (Open~5.)
\end{todo}

\begin{todo}
\semibold{The $\partial\Sigvol{}/\partial u$ decision.} Whether to retain
the state dependence of the noise amplitude in the response sector, and
if so whether to integrate the resulting Riccati flow directly or through
the symplectic block linearization, is to be decided empirically on the
gradient-free instanton, where $\Hfp$ is exactly conserved and the
closed form \eqref{eq:hfp-closed-form} isolates the contribution of these
terms. (Open~4; Section~\ref{sec:instanton-eqs}.)
\end{todo}

\begin{todo}
\semibold{The decoupled limit and the missing reduction test.} With the
gradient and advection terms switched off, the non-core response fields
vanish identically, but the core equations still carry the dilution
$\ncount(1,\Ny)=\tfrac32\Deltas$ and the response field must be written
$\rfield=\delta(\yco-1)\rfield_\core/\volel(1,\Ny)$ for the action to reduce
to the single-trajectory one. Whether the on-shell action values agree in
that limit was left as an empirical check, and the check was blocked: with
the core value eliminated through a Neumann row from background
neighbours, the zero row sum of $D$ pins $\field_\core$ to the background
for every $\lag$, so the fully decoupled shooting problem is degenerate.
No non-trivial reduction test exists. (\texttt{RECONSTRUCTION.md}, A.1.)
\end{todo}

\begin{todo}
\semibold{The MSR--spectral bridge.} The relation
$\Hfp=-\partial S/\partial\dNstar$ translates between the asymptotic
slope of the instanton action and a spectral gap $\Lambda_0$ only if both
are known independently; the claim that $\Hfp\to-\Lambda_0$ at large
$\dNstar$ was withdrawn on 12 July 2026. A Gaussian expansion about the
noiseless trajectory gives an action quadratic in $\dNstar$,
$S_{\min}=\dNstar^2/\sum_ib_i^2/\lambda_i$, and stalled on the
diagonalization step (the reduced weight is not the Sturm--Liouville
measure). The linear law $\Lambda_0\dNstar$ of the first-passage
formalism is a property of the marginalized object; the fixed-profile
instanton is one channel of that sum. The scaffold for both questions is
the adjoint backward-Kolmogorov eigenproblem, which has not been set up.
This is a research item. (Open~7; \texttt{RECONSTRUCTION.md}, B.2.)
\end{todo}

\begin{todo}
\semibold{Non-Markovian noise.} The local amplitude $\Sigvol{}\propto\Hub^2$
keeps the Langevin system Markovian and the phase space of $\Hfp$
unchanged. A diffusion model with memory of the prior state through the
mode functions (the Mukhanov--Sasaki equation that produces the
horizon-exit noise) is non-Markovian, puts the Freidlin--Wentzell
Hamiltonian on an extended phase space, and changes the canonical
structure of $\Hfp$, not only its coefficients. The present
construction is the Hamiltonian of the truncated Markovian problem.
(\texttt{RECONSTRUCTION.md}, B.2.)
\end{todo}

% ======================================================================
\section{Summary}
\label{sec:summary}
% ======================================================================

The gradient-corrected instanton is the saddle point of the action
\eqref{eq:msr-action} on the fixed rectangle $\yco\in[-1,1]$
(exterior to the core's own horizon), $\Ny\in[\Ninit,\Nfinal]$, with the
measure the physical volume element selected by the noise
(Section~\ref{sec:noise-measure}) and the inner edge anchored to the
core's horizon with its rate of change derived by the chain rule
(Section~\ref{sec:coordinate}). The saddle-point equations are Hamilton's
equations of the Fokker--Planck Hamiltonian $\Hfp$
(Section~\ref{sec:hfp}); the response fields are its conjugate momenta,
and the full response equations are generated by differentiating a single
discrete Hamiltonian rather than transcribed
(Section~\ref{sec:collocation-scheme}). At the core exactly one
characteristic enters in each sector, the Neumann condition on $\field$ is
the data-free reflecting closure, and no further condition on
$\mom(1,\Ny)$ may be imposed (Section~\ref{sec:bcs}). The problem is
closed by the terminal condition \eqref{eq:terminal-colloc} and the
shooting parameter $\lag$, and solved by Picard iteration. The corrected
core trajectory $\field(1,\Ny)$ and the profile $\zprofile(\yco)$ of
\eqref{eq:zeta-extraction} are outputs of the solution.

The numerical scheme described in Sections~\ref{sec:collocation-basis}
and~\ref{sec:collocation-scheme} is the target scheme; the single
characteristic penalty in the weighted norm and the differentiated
response sector are not yet implemented, and no action value computed
before October 2026 should be used, because the response sector then in
use omitted leading-order terms (Section~\ref{sec:instanton-eqs}).
The theoretical issues of Section~\ref{sec:open-issues} remain open: the
probability interpretation of the two-dimensional action, the largest-time
justification of the terminal condition, the response-sector indicial
analysis, the treatment of the state dependence of the noise amplitude,
the decoupled-limit reduction test, the bridge to the spectral formalism,
and the extension to non-Markovian noise. A discrete (finite-shell)
implementation (Section~\ref{sec:geometry}) remains available as an
independent cross-check on the collocation scheme.

\end{document}
